% Modernised reading — fonds Alexandre Grothendieck, shelfmark 77, pages 1-157
% (the whole folder).
%
% This is an INTERPRETATION, not a transcription. It re-reads the pages in
% today's notation and vocabulary, states the hypotheses the manuscript leaves
% implicit, and drops the critical apparatus so that the mathematics reads
% straight through. Everything the brackets and underlines carried moves into a
% footnote. It is held to being mathematically correct as it stands; where the
% manuscript is loose or mistaken, the true statement is what appears here and
% the footnote says what the page has. In any dispute of substance the
% transcription governs: whoever cites Grothendieck cites batch-01.fr.tex to
% batch-08.fr.tex (pp. 1-20, 21-40, 41-60, 61-80, 81-100, 101-120, 121-140,
% 141-157).
%
% Pass: Opus 5.5 (claude-opus-5-5), 2026-10-10 — first pass, unchecked against the pages by a human.
% The folder was read whole, all eight batches together; this is one document
% for the shelfmark. It was made on Opus 5.5 and has not been measured against
% the reference reading of folder 115 (made on Opus 5). Checked by machine
% (sympy) while writing: the 5x5 determinant of p. 12, factored as
% (alpha-2)^5 (alpha+2)^2 alpha (alpha+1)^2; the determinants of pp. 10 and 13
% and the tangent plane of p. 19; the identity of p. 5; the formula
% q([x,y]) = -4q(x)q(y) + phi(x,y)^2 of pp. 93 and 115; the boxed identity of
% p. 116 (wrong as written, corrected); the graph expansion of the universal
% discriminant (pp. 65-71, 87) and the congruence B^2 = (-1)^n Delta mod 4 for
% ranks 2, 4, 6, with Delta = 0 mod 2 in ranks 3, 5. The rest was checked by
% hand. The literature named in footnotes is cited from memory and was not
% checked against the sources.
%
% Scope. 157 facsimile pages, 145 transcribed. Pages 1, 6, 20, 64 and 126 are
% covers in his hand whose words head the parts here; pages 70, 72, 74, 76, 78,
% 80 and 82 are unrelated typed versos, not transcribed and not described.
% Page 88 is a pink folder carrying only a title.
%
% Spine. Quadratic forms and quadratic algebras over an arbitrary base, and
% what changes at the prime 2. Six parts, each under one of his covers:
%   pp. 1-5     I « Formes quadratiques de rg 2, droites projectives et
%               revêtements quadr. »: ten equivalent categories (forms of an
%               object with automorphism group G_m x| Z/2); reflections of a
%               conic and étale quadratic subalgebras of a quaternion algebra.
%   pp. 6-19    II « Formes quadratiques inv. et quasi-inv. sous D_infty,
%               D_n »: two drafts (pp. 7-12 + 17; pp. 18, 13-16, 19) of the
%               invariant conic circumscribed to a regular polygon.
%   pp. 20-63   III « Extensions quadratiques », his §§ 1-10: affine envelope,
%               splittings (b, c), discriminant, symmetric powers of an affine
%               line, algebra structures = quadratic forms with q(e) = 1 = norm,
%               canonical involution, 2 invertible: (L, D); char. 2:
%               Artin-Schreier; mixed case, 2 regular: (L, D, b_0) with
%               D_1 = gamma(b_0), in three successive versions.
%   pp. 64-87   IV « Discriminant des formes quadratiques »: the universal
%               discriminant as a sum over graphs; = 0 mod 2 in odd rank; in rank
%               2n, B^2 - (-1)^n Delta = 4C (B the sum over perfect matchings);
%               the ring version and the Proposition of p. 85; p. 87 the rank-4
%               check.
%   pp. 88-125  V « Quadriques affines »: § 1 (pp. 89-101) the affine quadric
%               Sigma: q = -1 of a special quadratic module of rank 3 is the
%               complement of the conic X in the Weil restriction X^{S[i]};
%               § 1' (pp. 102-116) the twisted version X^Omega ≃ Sigma-hat and
%               the sign check; § 1'' (pp. 117-125) « bi-regular » affine
%               quadrics, (T, q, a), canonical symmetry, automorphism group.
%   pp. 126-157 VI « Formes quadratiques. Fibrés en drapeaux quad., groupes
%               orthogonaux »: half-discriminant, regularity, O and SO in odd
%               and even rank, classification by special forms; isotropic flag
%               schemes Delta_H (smooth projective, connected fibres except one
%               case), the Arf double cover, stabilisers O_D, O_D^+, dimensions,
%               SO(q) smooth with connected fibres (in even rank defined as the
%               kernel of O(q) -> Z/2 given by the action on Arf).
%
% Notation fixed for the whole document.
%  - S a scheme; « Module » = locally free O_S-module of finite rank, identified
%    with its vector bundle. q a quadratic form, phi its polar form,
%    q(x+y) = q(x) + q(y) + phi(x,y). delta(q) = det of phi, a section of
%    (det E)^{-2}; delta'(q) the half-discriminant (his « discriminant
%    divisé »): delta(q) in even rank, delta(q)/2 in odd rank (well defined
%    integrally by p. 68).
%  - Part III: B a quadratic algebra, e its unit, L = (B/O_S)^vee (his
%    L-check = B/O_S), tr, N, sigma the canonical involution; D(B) = b^2 - 4c
%    in Gamma(L^2) the discriminant, written \mathfrak{D} throughout (his
%    \mathfrak{D}, Delta, round d, theta); E_Omega the affine envelope (his E,
%    pp. 21-27); F an affine line (L-torsor) in general (his F from p. 27);
%    the isomorphism det V ≃ L of p. 35 is kappa (his alpha). S_0 = V(2),
%    S_1 = V(4); gamma(b_0) the square of any lift.
%  - Part II: zeta, zeta^{-1} the eigenvalues of the rotation u,
%    alpha = zeta + zeta^{-1}.
%  - Part IV: Y_i, Y_ij the coefficients of the universal form; Delta its
%    discriminant, B the matching sum, Gamma_I the « graphs ».
%  - Part V: (E, q, omega) special quadratic module of rank 3 with
%    delta'(q) = -omega^{-2}; M the (twisted) quaternion algebra, gamma M its
%    elements of reduced trace 0, q_M = -det; S[i] = Spec O_S[T]/(T^2+1);
%    X the conic, X-tilde its punctured cone, Sigma the affine quadric
%    q = -1, X^{S[i]} the Weil restriction; Omega the étale double cover of
%    § 1'.
%  - Part VI: (E, q) regular of rank N (his n, then N); Delta_H the isotropic
%    flags of type H; Arf, Arf'; O_D, O_D^+, SO_D; R_i = E_i/E_{i-1},
%    Q = F^perp/F.
%  Homonymies: E is the rank-3 vector space of II, the affine envelope of III,
%  the rank-3 module of V § 1, the affine bundle of V § 1', the quadratic
%  module of VI; Delta is a degree-2 divisor (p. 2), the 5x5 determinant
%  (p. 12), the universal discriminant (IV); alpha is zeta + zeta^{-1} (II)
%  and, on p. 23 only, a trivialisation of L. Each is fixed where it occurs.
%
% Substantive departures from the pages, with page numbers:
%  - p. 2: « X conique lisse de X » read « de P ».
%  - p. 3: the margin « faux » is kept: T^{⊗2} is not Zariski-locally trivial.
%  - p. 7: on the dual, D_1-check^2 and D_2-check^2 have eigenvalues
%    zeta^{-2}, zeta^2 (page: zeta^2, zeta^{-2}); p. 8: V and W read
%    D_1-check^2 + D_2-check ⊗ D_0-check, etc.
%  - p. 10: zeta^4 - 1 invertible <=> alpha(alpha^2 - 4) invertible (page:
%    4(alpha^2 - 4)); the conclusion on non-degenerate forms holds as soon as
%    alpha^2 - 4 is invertible (the square case included, by p. 9); det of
%    the matrix of (5) is -2ab^2 (page 2ab^2), so delta' = -ab^2.
%  - p. 12: the determinant is (alpha-2)^5 (alpha+2)^2 alpha (alpha+1)^2 up to
%    sign (computed here; the page asks for the multiplicities).
%  - p. 18: « alpha - 1 inv. » read « 1 + alpha inv. »; « v = -b » read v = -c.
%  - p. 19: B = zeta - zeta^{-1} = -zeta^{-1}(1 - zeta^2) (page: +).
%  - p. 23: first matrix of mu_u has +bu (page -bu).
%  - p. 26: (24) and (13) agree in sign (the transcription's note says not).
%  - p. 33, (46): Sym^2(V) is an extension of O by L ⊗ V (page: by V).
%  - p. 41: (67) a) read g(e_B) = 1.
%  - p. 44: c_0 in Gamma(L^2) (page reads L^{-2}).
%  - pp. 45-46: the pseudo-reflection id - tr ⊗ e is -sigma; its « fixed »
%    line O·e is the line where it acts by -1.
%  - p. 49: B = Sym(L-check)/(w - (D/4)(w)) read for « (1 - 4 d) L^{-2} ».
%  - p. 56: a_1 does not depend on q_0; condition a) is implied by b).
%  - p. 57, Cor. 1: the density of U ∩ S_1 follows from that of U ∩ S_0.
%  - p. 58 b), p. 60, p. 62: the H^1 conditions made precise (F trivial iff
%    b_0 lifts to Gamma(S, L), automatic if 2 is injective on H^1(S, L)).
%  - p. 66: nu_* counts cycles of length >= 3, nu_*^+ those of even length >= 4.
%  - p. 68: B^2 = sum M + 2 sum over unordered pairs (page: ordered).
%  - p. 79: u in A_Delta^* (page A^*). p. 81: U = (2T + b)/u (page
%    (T + b/2)/2u).
%  - p. 85: proof of c) => a) supplied (localisation at 2A, the margin idea of
%    p. 83); uniqueness of B_Delta stated by the page, not proved.
%  - p. 89, (1): delta'(q) = -omega^{⊗-2} (page delta(q) = -omega^{⊗2}).
%  - p. 92: q(ax - by) has -ab phi(x,y) (page -2ab).
%  - p. 98: Sigma ≃ X^{S[i]} - X (page: with tildes).
%  - p. 100: lambda^{-3} mu^{-2} = -2^{-2} (page -2^2).
%  - p. 108: (E, -q, omega) is special for q_M = -det, i.e. q = det on gamma M.
%  - pp. 114-115: F_M(alpha, beta) = tr(alpha beta) - [alpha, beta] = 2 beta
%    alpha vanishes on the diagonal; it gives the isomorphism off the
%    diagonal only.
%  - p. 116, (3): det(u - s) = det u - s^2; the boxed identity has
%    + det u (tr v)^2 + det v (tr u)^2 (page: minus signs).
%  - p. 120 N.B. a): with pi(a) = 2, the equation is q_0(x) = -1/4 before
%    rescaling (page q_0(x) = -1).
%  - p. 123: sigma^2 = id because pi(a) = 2 (page « pi(a) = 1 »); « rg E
%    pair » read T of even rank.
%  - p. 124: polar form of det is ad' + da' - bc' - cb' (page: opposite), so
%    pi = Tr.
%  - p. 128: B_nu is simple for all nu >= 1 (page « simple si n >= 7 »).
%  - p. 131, Cor: Aut(X) ≃ O(q)/mu_2 for N >= 3.
%  - p. 138: e_1, ..., e_n a basis of F (page: of E). p. 139: delta(q) =
%    (-1)^n omega_F^{-2} (page (omega_F)^2).
%  - p. 141: F' ⊂ F (page F ⊂ F').
%  - p. 148: last term of the completed flag is E (page F).
%  - p. 149, Cor. 2: étale-local only if 2 is invertible or rank Q is even;
%    fppf-local in general.
%  - p. 152: the map u -> q(y + u y) is not linear; its zero locus is still,
%    over affine S, isomorphic as a scheme to a vector bundle of the stated
%    dimension.
%  - p. 154: Q = 0 (page « de rg 1 »); the first definition SO = O ∩ SL in
%    even rank fails in char. 2, as p. 155 says.
%
% Announced and never established in the folder: the equivalence of the ten
% categories of p. 2 (asserted); what the « invariants » of the regular
% spherical polyhedra of p. 89 are (the folder sets up Sigma and stops); the
% p. 14-16 identification of the conic with an affine line and the group
% scheme of p. 16 (struck); the global extension problem of p. 86; the
% structure G ≃ G^0 x Z/2 of p. 125 (« on va montrer »); the char. 2 attempt
% of p. 99 and the rescaled formula (26) of p. 100; the transitivity
% statement in the margin of p. 155 (no proof).
\documentclass[11pt,a4paper]{article}
\input{../preamble/grothendieck}

\folder{77}
\pages{1}{157}
\dating{[à partir de 1977]}
\watermark{Édition de démonstration\\interprétation personnelle de l'œuvre}
\foldertitle{Quadriques affines, extensions quadratiques : notes manuscrites (s.d.) — lecture modernisée du dossier entier}

\begin{document}

\section*{Résumé}

\begin{resume}
Ces pages parlent des formes quadratiques — les expressions comme $x^2 + y^2$,
$xy$ ou $x^2 + bxy + cy^2$ — et de ce qu'elles deviennent quand on cesse de
travailler sur les nombres réels ou complexes pour travailler « sur une base
quelconque » : sur les entiers, sur un anneau où $2 = 0$, ou sur une famille
d'anneaux qui varie avec un paramètre.

Le problème est ancien et simple à dire. Sur les réels, une équation
$x^2 + bx + c = 0$ se résout en divisant par $2$ : on complète le carré, et
tout est gouverné par le discriminant $b^2 - 4c$. Sur les entiers, cette
division est interdite, et quelque chose se perd : deux équations de même
discriminant peuvent cesser d'être équivalentes, et un discriminant n'est pas
n'importe quel nombre — il est toujours congru à $0$ ou à $1$ modulo $4$. Dans
un anneau où $2 = 0$, le discriminant n'est plus qu'un carré et ne dit presque
plus rien. Tout le dossier est une manière de tenir les comptes de ce qui se
passe au-dessus du nombre premier $2$, au lieu de l'exclure d'emblée.

L'idée qui revient partout est géométrique. Une équation du second degré a
deux racines ; la donnée de ses racines, prises ensemble, est un
« revêtement à deux feuillets » de la base, comme les deux déterminations de
$\sqrt{t}$ au-dessus du plan complexe. À toute forme quadratique,
Grothendieck attache ainsi, de plusieurs façons, un tel revêtement double :
par le discriminant quand $2$ est inversible, par une équation d'Artin et
Schreier quand $2 = 0$, et, dans le cas mixte des entiers, par un discriminant
muni d'une racine carrée modulo $2$ qui se prolonge modulo $4$. Un revêtement
double du même genre reparaît à la fin sous un autre nom : les sous-espaces
isotropes maximaux d'une forme de rang pair forment deux familles, et le choix
d'une famille est un revêtement double de la base — celui qui porte
l'invariant d'Arf.

Le dossier commence par la géométrie la plus concrète : la conique qui passe
par les sommets d'un polygone régulier, et les réflexions d'une droite
projective. Il passe ensuite à l'algèbre — les « extensions quadratiques »,
classées d'abord quand $2$ est inversible, puis en caractéristique $2$, puis
dans le cas mixte —, à un calcul combinatoire du discriminant d'une forme
générale par des graphes, à la sphère vue comme une quadrique affine, et
s'achève sur les variétés de drapeaux isotropes et les groupes orthogonaux
sur une base quelconque, la caractéristique $2$ comprise.

Une analogie aide à lire la partie sur la sphère. La sphère réelle
$x^2 + y^2 + z^2 = 1$ et la droite projective complexe sont la même chose :
c'est la sphère de Riemann. Un point de la sphère, vu de cette façon, est un
couple de points « conjugués » de la droite projective complexe. Grothendieck
écrit cette identification sans les nombres réels ni complexes, avec
seulement l'équation $T^2 + 1 = 0$ adjointe à la base, et elle devient un
énoncé valable sur tout anneau où $2$ est inversible.

Les noms modernes à chercher ensuite : algèbres quadratiques et leur
discriminant, congruence de Stickelberger, revêtements d'Artin--Schreier,
restriction des scalaires de Weil, isomorphisme exceptionnel
$\mathrm{PGL}_2 \simeq \mathrm{SO}_3$, grassmanniennes orthogonales, invariant
d'Arf et invariant de Dickson, sous-groupes paraboliques des groupes
orthogonaux.

\keywords{quadratic algebra, discriminant of a quadratic algebra, quadratic étale cover, affine hull, symmetric power of an affine line, Hilbert scheme of points, canonical involution, Artin-Schreier cover, Stickelberger congruence, half-discriminant, hafnian, Pfaffian, signed discriminant, dihedral group, circumscribed conic, regular polygon, quaternion algebra, Weil restriction, affine quadric, exceptional isomorphism PGL2 = SO3, Segre embedding, normalizer of a maximal torus, orthogonal Grassmannian, isotropic flag variety, Stein factorization, Arf invariant, Dickson invariant, special orthogonal group scheme, parabolic subgroup}
\end{resume}

\pagerange{2}{157}
\section*{Le fil du dossier, et les conventions}

Le dossier est fait de six liasses, chacune sous une couverture de sa main
dont les mots servent ici de titres de parties. Les nombres au crayon de
certaines couvertures sont des comptes de pages et tombent juste :
« (13~p) » sur la couverture de la page 6 couvre les pages 7 à 19 ;
« (43~p) » sur celle de la page 20 couvre les pages 21 à 63 ; « (31~p) » sur
celle de la page 126 couvre les pages 127 à 157 ; « (16p) » sur celle de la
page 64, qui n'est peut-être pas de lui, tombe juste aussi : les pages 65 à 87
en comptent seize une fois ôtés les sept versos dactylographiés. Les numéros de formules
sont ceux des pages ; ils repartent à chaque liasse, parfois à chaque
rédaction.

\subsection*{Les stations}

\begin{itemize}
  \item \textbf{I. Formes quadratiques de rang $2$, droites projectives et
    revêtements quadratiques (pages 1 à 5).} Dix catégories équivalentes,
    toutes formes d'un même objet dont le groupe d'automorphismes est
    $\mathbb{G}_m \rtimes \mathbf{Z}/2$ ; puis les réflexions d'une conique,
    les sous-algèbres étales de rang $2$ d'une algèbre de quaternions, et la
    formule qui fait d'une involution $\sigma$ une réflexion orthogonale.
  \item \textbf{II. Formes quadratiques invariantes et quasi-invariantes sous
    le groupe diédral (pages 6 à 19).} Deux rédactions du même calcul : la
    conique invariante circonscrite à un polygone régulier, d'abord en
    coordonnées diagonales (pages 7 à 12 et 17), puis en « coordonnées
    universelles » (pages 18, 13 à 16 et 19).
  \item \textbf{III. Extensions quadratiques (pages 20 à 63)}, ses §§ 1 à 10 :
    l'enveloppe affine d'un revêtement quadratique ; les scindages et les
    sections $b$, $c$ ; le discriminant ; les puissances symétriques d'une
    droite affine ; le théorème « structures d'algèbre $=$ diviseurs de degré
    $2$ $=$ formes quadratiques valant $1$ en $e$ », la forme étant la norme ;
    l'involution canonique ; la classification quand $2$ est inversible, en
    caractéristique $2$ (Artin--Schreier), et dans le cas mixte où $2$ est
    seulement régulier, en trois rédactions successives (pages 52 à 57, 58 à
    60, 61 à 63).
  \item \textbf{IV. Discriminant des formes quadratiques (pages 64 à 87).} Le
    déterminant de la forme quadratique universelle comme somme sur des
    graphes ; sa nullité modulo $2$ en rang impair ; en rang $2n$, l'égalité
    $B^2 - (-1)^n \Delta = 4C$ ; la version pour un anneau et un élément
    $\Delta$, jusqu'à la Proposition de la page 85 ; la vérification en rang
    $4$ (page 87).
  \item \textbf{V. Quadriques affines (pages 88 à 125).} § 1 : la quadrique
    affine $q = -1$ d'un module spécial quadratique de rang $3$ est le
    complémentaire de la conique $X$ dans la restriction de Weil
    $X^{S[\mathbf{i}]}$ ; § 1' : la version tordue $X^{\Omega} \simeq
    \hat{\Sigma}$ et le contrôle du signe ; § 1'' : les quadriques affines
    « birégulières », leur symétrie canonique et leur groupe d'automorphismes.
  \item \textbf{VI. Formes quadratiques, fibrés en drapeaux, groupes
    orthogonaux (pages 126 à 157).} Discriminant divisé et régularité ; $O$ et
    $\mathrm{SO}$ en rang impair et en rang pair ; les schémas de drapeaux
    isotropes, lisses et projectifs, à fibres connexes sauf en un cas ; le
    revêtement d'Arf ; les stabilisateurs et leurs dimensions ; $\mathrm{SO}(q)$
    lisse à fibres connexes, défini en rang pair par l'action sur le
    revêtement d'Arf.
\end{itemize}

\subsection*{Conventions, valables pour tout le dossier}

$S$ est un schéma ; « Module » veut dire $\mathcal{O}_S$-module localement libre
de rang fini, et un fibré vectoriel est identifié au Module de ses sections.
Une forme quadratique $q$ et sa forme polaire $\varphi$ sont liées par
$q(x + y) = q(x) + q(y) + \varphi(x, y)$, d'où $\varphi(x, x) = 2q(x)$. Le
\emph{discriminant} $\delta(q)$ est le déterminant de $\varphi$, section de
$(\det E)^{\otimes -2}$ ; le \emph{discriminant divisé} $\delta'(q)$ est
$\delta(q)$ en rang pair et $\delta(q)/2$ en rang impair — division légitime
sur toute base, puisque le déterminant de la forme universelle de rang impair
est divisible par $2$ (page 68). C'est ce qu'on appelle aujourd'hui le
demi-discriminant. Une forme est \emph{régulière} (ou lisse) si $\delta'(q)$
est inversible.

Dans la partie III, $B$ est une algèbre quadratique (commutative, localement
libre de rang $2$), $e = e_B$ son unité, $\mathrm{tr}$ et $N$ sa trace et sa
norme, $\sigma$ son involution canonique, et
\[
  L = (B/\mathcal{O}_S)^{\vee},
\]
de sorte que le quotient $B/\mathcal{O}_S$ est $L^{\vee}$ — la page écrit
$\check{L} = B/\mathcal{O}_S$. Le discriminant d'une algèbre quadratique est
noté $\mathfrak{D}$ partout, bien que les pages écrivent tour à tour
$\mathfrak{D}$, $\Delta$, $\partial$ et $\vartheta$ ; c'est une section de
$L^{\otimes 2}$. On pose $S_0 = V(2\mathcal{O}_S)$ et $S_1 = V(4\mathcal{O}_S)$.

\textbf{Homonymies.} La lettre $E$ désigne successivement l'espace de dimension
$3$ de la partie II, l'enveloppe affine d'un revêtement quadratique (partie
III, pages 21 à 27, où on l'écrit $E_{\Omega}$), le module de rang $3$ de la
partie V, § 1, le fibré affine de rang $3$ du § 1' et le module quadratique de
la partie VI ; la droite affine générale de la partie III est $F$, comme sur les
pages à partir de la page 27. La lettre $\Delta$ désigne un diviseur de degré
$2$ (page 2), un déterminant $5 \times 5$ (page 12), le discriminant universel
(partie IV) et, dans la partie VI, les schémas de drapeaux $\underline{\Delta}_H$
— on garde la notation des pages, le contexte ne laissant jamais de doute.
Enfin $\alpha = \zeta + \zeta^{-1}$ dans la partie II ; l'isomorphisme
$\det V \simeq L$ que la page 35 appelle aussi $\alpha$ est écrit $\kappa$ ici.

\subsection*{Ce que seule une lecture d'ensemble peut dire}

\begin{itemize}
  \item \emph{Le revêtement double d'une forme quadratique est le fil du
    dossier.} Il apparaît comme l'algèbre étale $A$ de la catégorie $C_2$
    (page 2), comme le revêtement quadratique de discriminant $\mathfrak{D}$
    quand $2$ est inversible (page 48), comme le revêtement d'Artin--Schreier
    en caractéristique $2$ (page 51), comme le revêtement $\Omega$ des racines
    carrées de $\delta'(q)$ (page 105), et comme le schéma $\underline{\mathrm{Arf}}$ des
    composantes connexes de la variété des sous-espaces isotropes maximaux
    (page 142). La page 144 relie les deux dernières par
    $\underline{\mathrm{Arf}}' \simeq \underline{\mathrm{Arf}} \wedge^{\mathbf{Z}/2} \mu_2$.
  \item \emph{Le « cas mixte » et le discriminant universel sont la même
    congruence.} La condition $\mathfrak{D}_1 = \gamma(b_0)$ du théorème des
    pages 58 à 60 (une racine carrée de $\mathfrak{D}$ modulo $2$, dont le carré
    se relève modulo $4$) est ce que le Lemme de la page 73 isole pour un
    anneau, et ce que la page 68 établit pour la forme universelle de rang pair
    : $(-1)^n \Delta \equiv B^2 \pmod 4$. À coordonnées fixées, le couple
    $(B, C)$ de la page 68 est donc, par le § 10, un revêtement quadratique de
    discriminant $(-1)^n \delta(q)$. Que ce revêtement ne dépende pas des
    coordonnées, et qu'il soit le revêtement d'Arf de la page 142, les pages
    ne le disent pas\footnote{Rapprochement de l'édition. Aujourd'hui ce
    revêtement est le spectre du centre de l'algèbre de Clifford paire, et la
    congruence est une forme de celle de Stickelberger ; le dossier 84 (pages
    48 à 53) énonce le même théorème de classification avec les triplets
    $(L, \delta, T_0)$, et le rapproche lui aussi de Stickelberger.}.
  \item \emph{La quadrique affine de la partie V est le schéma des
    involutions $\sigma$ de la page 5.} Les éléments $\sigma$ d'une algèbre de
    quaternions, de trace nulle et de déterminant $-1$ (donc $\sigma^2 = 1$),
    sont les points de la quadrique $\Sigma$ des pages 108 et 109 ; la page 5
    les met en bijection avec les diviseurs étales de degré $2$ de la conique,
    les pages 89 à 116 avec les couples ordonnés de points distincts de la
    conique. Les deux énoncés se correspondent : $\sigma$ et $-\sigma$ donnent
    la même réflexion et les deux ordres d'un même couple.
  \item \emph{Les trois réflexions de la page 4 sont celles du groupe
    cartographique.} Les hypothèses A) à D) — $\tau_0\tau_2 = \tau_2\tau_0$,
    $\tau_0\tau_1$ et $\tau_1\tau_2$ non commutatives — sont celles que les
    dossiers 86 et 88 imposent à $\sigma_0, \sigma_1, \sigma_2$ ; le titre
    de la page 89, « Invariants des groupes des polyèdres réguliers
    sphériques », dit où cela devait mener.
\end{itemize}

\subsection*{Ce que le dossier annonce sans l'établir}

L'équivalence des dix catégories de la page 2, affirmée ; les invariants des
groupes de polyèdres réguliers que promet la page 89 — le dossier construit la
quadrique et s'arrête ; l'identification de la conique à une droite affine et
le schéma en groupes des pages 14 à 16 (biffés) ; le problème global de
prolongement de la page 86 ; la décomposition $G \simeq G^0 \times \mathbf{Z}/2$ de
la page 125 (« on va montrer ») ; la tentative en caractéristique $2$ de la
page 99 et la formule (26) de la page 100, interrompue ; l'énoncé de
transitivité en marge de la page 155, sans démonstration.

\section*{I. Formes quadratiques de rang $2$, droites projectives et revêtements quadratiques (pages 1 à 5)}

\pagerange{2}{3}
\subsection*{Dix catégories équivalentes (pages 2 et 3)}

Sous le titre « Digression quadratique », la page 2 aligne, pour un schéma $S$,
dix catégories (on n'y considère que les isomorphismes) :
\begin{itemize}
  \item $C_1(S)$ : les Modules quadratiques réguliers $(V, q)$ de rang $2$ ;
  \item $C_2(S)$ : les triples $(A, V, \xi)$, où $A$ est une $\mathcal{O}_S$-algèbre
    finie étale de rang $2$, $V$ un $A$-Module inversible et $\xi$ une
    trivialisation de $N_{A/S}(V)$ ;
  \item $C_3(S)$ : les couples $(U, T)$, où $U$ est une forme de $\mathbb{G}_m$ (un
    tore de rang $1$) et $T$ un $U$-torseur ;
  \item $C_4(S)$ : les couples $(X, \Delta)$, où $X$ est un fibré en droites
    projectives et $\Delta \subset X$ un diviseur relatif étale de degré $2$ ;
  \item $C_5(S)$ : les triples $(P, X, D)$, où $P$ est un fibré en plans
    projectifs, $X \subset P$ une conique lisse et $D \subset P$ un sous-fibré en
    droites, non tangent à $X$ en chaque fibre\footnote{La page écrit
    « $X$ conique lisse de $X$ » ; on lit « de $P$ ».} ;
  \item $C_6(S)$ : les couples $(E, C)$, où $E$ est un fibré affine de rang $2$
    et $C \subset E$ une conique projectivement lisse, non parabolique en chaque
    fibre ;
  \item $C_7(S)$ : les quadruples $(E, Q, \omega, V)$, où $(E, Q, \omega)$ est un
    module spécial orthogonal de rang $3$ et $V \subset E$ un sous-fibré de rang
    $2$ non isotrope, c'est-à-dire tel que $Q|V$ soit régulière ;
  \item $C_8(S)$ : les couples $(M, A)$, où $M$ est une algèbre d'Azumaya de rang
    $4$ et $A \subset M$ une sous-algèbre commutative étale de rang $2$ ;
  \item $C_9(S)$ : les formes $G$ de $\mathrm{PGL}_{2,S}$ munies d'un tore maximal,
    et $C'_9(S)$ : les formes de $\mathrm{SL}_{2,S}$ munies d'un tore
    maximal\footnote{Écrits en marge, la feuille tournée ; la page note
    $\mathrm{GP}(1)$ pour $\mathrm{PGL}_2$, comme le dossier 82.} ;
  \item $C_{10}(S)$ : les torseurs sous $\mathbb{G}_{m,S} \rtimes (\mathbf{Z}/2)_S$, le
    produit semi-direct où $\mathbf{Z}/2$ opère par $t \mapsto t^{-1}$\footnote{La
    page écrit le produit semi-direct avec un signe propre, entre les deux
    facteurs, que la transcription rend par « $\tfrac{1}{2}$ ».}.
\end{itemize}
« Toutes ces catégories sont canoniquement équivalentes… » La page ne le
démontre pas. La raison est que chacune est la catégorie des formes (tordues)
d'un objet type défini sur $\mathbf{Z}$ dont le schéma en groupes
d'automorphismes est $\mathbb{G}_m \rtimes \mathbf{Z}/2$ : le groupe orthogonal de
$xy$, le normalisateur du tore diagonal dans $\mathrm{PGL}_2$, le stabilisateur
de $\{0, \infty\}$ dans $\mathrm{Aut}(\mathbb{P}^1)$, celui d'un plan non
isotrope dans $\mathrm{SO}_3$\footnote{Cette justification est de l'édition.
C'est le dictionnaire classique entre formes binaires, algèbres quadratiques
munies d'un module, tores de rang $1$ et paires « conique, diviseur de degré
$2$ » ; voir par exemple M.-A. Knus, \emph{Quadratic and Hermitian Forms over
Rings} (1991), cité de mémoire. Pour $C_1$, « régulier » doit s'entendre au sens
où l'algèbre de Clifford paire est étale, ce qui garde la caractéristique $2$.}.
Le passage de $C_1$ à $C_2$ est celui de l'algèbre de Clifford paire ; de $C_2$
à $C_3$, celui qui prend pour $U$ le tore des unités de norme $1$ de $A$ et
pour $T$ le torseur des isomorphismes de $(A, 1)$ sur $(V, \xi)$.

La page 3 cherche ensuite à rendre explicite la donnée d'une trivialisation du
torseur $T^{\wedge 2}$ déduit de $T$ par l'élévation au carré. Elle avance
d'abord qu'un tel torseur est localement trivial pour la topologie de Zariski,
puis se corrige en marge : « faux », comme le montrent sur $\mathbf{R}$ les formes
$\alpha(x^2 + y^2)$ et $xy$\footnote{Sur un corps, localement trivial pour
Zariski veut dire trivial ; or les formes $x^2 + y^2$ et $-(x^2 + y^2)$ donnent
sur $\mathbf{R}$ des torseurs distincts sous $U = \mathrm{SO}(2)$. La note marginale
est en partie illisible.}. Par la suite exacte de Kummer
\[
  0 \to \mu_2 \to U \xrightarrow{\ 2\ } U \to 0,
\]
trivialiser $T^{\wedge 2}$ revient à réduire le groupe structural à $\mu_2$,
c'est-à-dire à se donner une orbite de $T$ sous $\mu_2$ ; géométriquement, un
revêtement de degré $2$ de fibrés en droites projectives $X \to X'$, étale hors
des diviseurs $\Delta \subset X$ et $\Delta' \subset X'$, avec une section de
$X' \smallsetminus \Delta'$. Le groupe structural se réduit alors à
$\mathbf{Z}/2 \times \mu_2$, « le $4$-groupe », qui opère sur
$\mathcal{O}_S \times \mathcal{O}_S$ par symétrie et homothéties, et la catégorie
obtenue est celle des triples $(A, L, \varepsilon)$ : $A$ étale de rang $2$, $L$
inversible sur $S$, $\varepsilon$ une trivialisation de $L^{\otimes 2}$ ; on
reprend $V = A \otimes L$, d'où $N_{A/S}(V) \simeq L^{\otimes 2}$, et $\xi =
\varepsilon$\footnote{Deux notes marginales sur l'action de $\mu_2$ sur $(X,
\Delta)$ et sur $X' = X/\mu_2$ ne se lisent que par fragments ; elles ne sont
pas utilisées.}.

\pagerange{4}{5}
\subsection*{Réflexions d'une conique, sous-algèbres étales et involutions (pages 4 et 5)}

Soit $M$ une algèbre de quaternions (algèbre d'Azumaya de rang $4$) sur $S$,
$\gamma M$ le Module de ses éléments de trace réduite nulle, $C =
\check{\mathbb{P}}(\gamma M)$ le plan projectif correspondant et $X \subset C$ la
conique $\det = 0$ ; $X$ est un fibré en droites projectives. La page 5 aligne
une chaîne de bijections, fonctorielles en $S$ :
\begin{itemize}
  \item les réflexions de $X$ (involutions distinctes de l'identité sur chaque
    fibre) ;
  \item les réflexions du plan $C$ qui stabilisent $X$, ou les réflexions
    orthogonales de $\gamma M$ ;
  \item les points de $C \smallsetminus X$ (le centre de la réflexion) ;
  \item les droites non isotropes de $\gamma M$ ;
  \item les diviseurs étales de degré $2$ de $X$ (les points fixes de la
    réflexion) ;
  \item les droites de $C$ non tangentes à $X$ (la polaire du centre) ;
  \item les plans de $\gamma M$ non tangents au cône $q = 0$ ;
  \item les sous-algèbres étales de rang $2$ de $M$ ;
  \item les éléments $\sigma \in M$ tels que $\sigma^2 = 1$ et $\sigma \neq \pm 1$
    sur chaque fibre, c'est-à-dire de trace nulle et de déterminant $-1$.
\end{itemize}
Une accolade en marge réserve plusieurs de ces flèches à la caractéristique
$\neq 2$, et c'est nécessaire : en caractéristique $2$ un élément $\sigma$ de
trace nulle vérifiant $\sigma^2 = 1$ ne définit pas une algèbre
étale\footnote{La dernière condition est surchargée sur la page ; on lit
« $\det \sigma = -1$ » en remplacement d'une première version.}.

La réflexion $\tau$ de $X$ attachée à $\sigma$ agit sur $M$, par transport de
structure, par $u \mapsto \sigma u \sigma^{-1}$ ; « mais ce n'est pas une
réflexion de $M$ ! » — elle a les valeurs propres $1$ et $-1$, chacune de
multiplicité $2$. Pour obtenir une réflexion, il faut la composer avec
l'involution canonique $u \mapsto u^{*} = \mathrm{Tr}(u)\cdot 1 - u$, et l'on a
\[
  \tilde{\sigma}(u) = \sigma u^{*} \sigma = (\sigma u \sigma)^{*}
  = u - \mathrm{Tr}(u\sigma)\, \sigma .
\]
C'est exactement la réflexion orthogonale de $(M, \det)$ le long de $\sigma$ :
pour la forme $\det$, de polaire $\varphi(u, v) = \mathrm{Tr}\,u\,\mathrm{Tr}\,v -
\mathrm{Tr}(uv)$, on a $\varphi(u, \sigma) = -\mathrm{Tr}(u\sigma)$ et
$\det\sigma = -1$\footnote{Formule encadrée sur la page ; l'identité a été
vérifiée pour une matrice $2 \times 2$ générale $u$ et $\sigma$ de trace nulle
avec $\sigma^2 = 1$. L'interprétation comme réflexion orthogonale est de
l'édition.}.

La page 4, plus ancienne dans l'ordre de l'écriture semble-t-il, considère
trois réflexions $\tau_0, \tau_1, \tau_2$ de $X$, de centres $h_0, h_1, h_2 \in C
\smallsetminus X$, et pose les hypothèses suivantes :
\begin{itemize}
  \item[A)] $\tau_0 \tau_2 = \tau_2 \tau_0$, c'est-à-dire $h_0 \perp h_2$ ;
  \item[B)] $\tau_0\tau_1 \neq \tau_1\tau_0$ et $\tau_1\tau_2 \neq \tau_2\tau_1$ sur
    chaque fibre, c'est-à-dire $h_0 \not\perp h_1$ et $h_1 \not\perp h_2$ ;
  \item[C)] $h_0, h_1, h_2$ engendrent $C$ en chaque fibre ;
  \item[D)] la droite de $C$ formée des points fixes de $\tau_1$ et de $\tau_2$
    n'est tangente à $X$ en aucune fibre ; aux points de caractéristique
    $p > 0$ impaire, l'ordre de $\tau_1\tau_2$ doit être différent de $p$.
\end{itemize}
Les $\tau_i$ se relèvent de façon unique, au signe près, en des réflexions de
$M^{\times}$\footnote{La fin de D), au bas de la page, est très raturée ; seule
la condition sur l'ordre de $\tau_1\tau_2$ en caractéristique $p$ se lit avec
quelque assurance. Les notes marginales, reliées à B) et C), ajoutent que
$\tau_i^2 = 1$ avec $\tau_i \neq 1$ sur chaque fibre fait des $\tau_i$ des
réflexions, et demandent comment exprimer $h_i \notin X$ en termes des
$\tau_i$.}. Ce sont, pour la conique, les trois involutions génératrices du
groupe cartographique des dossiers 86 et 88, avec la seule relation
$(\tau_0\tau_2)^2 = 1$ : le point de départ d'une réalisation des polyèdres
réguliers comme groupes de réflexions de la sphère.

\section*{II. Formes quadratiques invariantes et quasi-invariantes sous le groupe diédral (pages 6 à 19)}

La couverture de la page 6 porte « Formes quadratiques inv.\ et quasi-inv.\ sous
$D_\infty$, $D_n$ » et « (13~p) ». Le cadre est celui du dossier 78 : un polygone
régulier $P$ dans un plan affine, d'invariant numérique $\alpha = \zeta +
\zeta^{-1}$, où $\zeta, \zeta^{-1}$ sont les valeurs propres de la rotation $u$
d'un pas ; le groupe diédral $\mathbb{D}_\infty$ (ou $\mathbb{D}_n$) engendré par
deux symétries $\sigma, \sigma'$ avec $\sigma' = u\sigma$ opère sur le plan, donc
sur l'espace vectoriel $E$ de dimension $3$ dont le plan est l'hyperplan
affine. Une forme quadratique est \emph{quasi-invariante} si elle est invariante
à un scalaire près, c'est-à-dire si elle engendre une droite invariante de
$\mathrm{Sym}^2(\check{E})$. Les pages 7 à 19 contiennent deux rédactions du même
calcul, dans deux systèmes de coordonnées ; l'ordre des feuillets les mêle.

\pagerange{7}{12}
\subsection*{Première rédaction : coordonnées diagonales (pages 7 à 12)}

\textbf{Cas I : $\alpha^2 - 4$ inversible.} Après le changement de base
quadratique $S' = \mathrm{Spec}\,\mathcal{O}_S[T]/(T^2 - \alpha T + 1)$, qui adjoint
$\zeta$, la représentation se diagonalise :
\[
  \sigma(x, y) = (y, x), \qquad \sigma'(x, y) = (\zeta y, \zeta^{-1} x), \qquad
  u(x, y) = (\zeta x, \zeta^{-1} y),
\]
d'où $E = D_1 \oplus D_2 \oplus D_0$, les droites propres de $u$ pour $\zeta$,
$\zeta^{-1}$, $1$. Sur $\mathrm{Sym}^2(\check{E})$, les six droites
$\check{D}_1^{2}, \check{D}_2^{2}, \check{D}_0^{2}, \check{D}_1\check{D}_2,
\check{D}_1\check{D}_0, \check{D}_2\check{D}_0$ sont propres pour $u$, de valeurs
propres $\zeta^{-2}, \zeta^{2}, 1, 1, \zeta^{-1}, \zeta$\footnote{La page écrit
$\zeta^{2}, \zeta^{-2}$ pour les deux premières ; sur le dual on a
$\zeta^{-2}, \zeta^{2}$. Rien n'en dépend.}.

Si les cinq valeurs $1, \zeta^{\pm 1}, \zeta^{\pm 2}$ sont distinctes sur chaque
fibre — c'est-à-dire si $\zeta^3 \neq 1$ et $\zeta^4 \neq 1$ en plus de $\zeta \neq
\pm 1$ —, les droites propres de $u$ sont les quatre droites $\check{D}_1^{2},
\check{D}_2^{2}, \check{D}_1\check{D}_0, \check{D}_2\check{D}_0$, que $\sigma$
échange deux à deux, et les droites de $\check{D}_0^{2} \oplus
\check{D}_1\check{D}_2$, sur lequel $\sigma$ opère trivialement. Les seules droites
invariantes par $\mathbb{D}_\infty$ sont donc formées de formes \emph{invariantes}.
La conclusion subsiste si l'on permet $\zeta^3 = 1$ ($\alpha = -1$, le triangle)
sur certaines fibres : $\sigma$ échange alors les deux plans propres
$\check{D}_1^{2} \oplus \check{D}_2\check{D}_0$ et $\check{D}_2^{2} \oplus
\check{D}_1\check{D}_0$\footnote{La page note ces plans « $\check{D}_1^{\otimes 2}
\oplus \check{D}_2$ » et « $\check{D}_2^{\otimes 2} \oplus \check{D}_1$ » ; on
complète le facteur $\check{D}_0$.}.

Le cas du carré, $\zeta^4 = 1$, c'est-à-dire $\zeta^2 = -1$ ou $\alpha = 0$ (alors
$2$ est inversible, puisque $\zeta \neq \pm 1$), fait apparaître deux droites
quasi-invariantes non invariantes : $\sigma$ permute $\check{D}_1^{2}$ et
$\check{D}_2^{2}$, et le plan qu'elles engendrent contient deux droites propres
pour $\sigma$, de valeurs propres $\pm 1$. Dans les coordonnées où les sommets du
carré sont $(\pm 1, \pm 1)$, ce sont les formes $x^2 - y^2$ et $xy$ ; l'une prend
deux valeurs non nulles sur les sommets, l'autre la valeur $+1$ sur deux sommets
opposés et $-1$ sur les deux autres, et la seule conique circonscrite qu'elles
fournissent, $x^2 - y^2 = 0$, est dégénérée : la réunion des deux diagonales.

\textbf{Conclusion} (page 10). Sur chaque fibre où $\alpha^2 - 4$ est inversible,
les formes quadratiques sur $E$ quasi-invariantes par $\mathbb{D}$ et non
dégénérées sont invariantes ; ce sont, dans les coordonnées diagonales, les
formes
\[
  (5) \qquad Q(x, y, z) = a z^2 + b\, xy ,
\]
de déterminant $-2ab^2$ et de discriminant divisé $-ab^2$ ; $Q$ est régulière si
et seulement si $a$ et $b$ sont inversibles\footnote{La page conclut sous
l'hypothèse « $\zeta^4 - 1$ inversible, i.e.\ $4(\alpha^2 - 4)$ inversible » :
$\zeta^4 - 1$ est inversible si et seulement si $\alpha(\alpha^2 - 4)$ l'est, et le
cas du carré montre que l'hypothèse $\alpha^2 - 4$ inversible suffit pour les
formes non dégénérées, les formes quasi-invariantes supplémentaires du carré
étant dégénérées. La page donne pour déterminant $2ab^2$ et pour discriminant
divisé $ab^2$ ; le déterminant de la matrice qu'elle écrit est $-2ab^2$.}. Les
sommets de $P$ sont les points $(\zeta^i, \zeta^{-i}, 1)$ ; $Q$ s'y annule si et
seulement si $a + b = 0$. Il existe donc une unique conique circonscrite à $P$
invariante sous $\mathbb{D}_\infty$,
\[
  (9) \qquad z^2 - xy = 0,
\]
et, en coordonnées affines, un unique polynôme de degré $\leq 2$ invariant et nul
sur les sommets, $xy - 1$ — ou encore une unique forme quadratique sur l'espace
des translations, invariante et égale à $1$ sur les sommets, $xy$.

\textbf{Cinq points suffisent} (pages 11 et 12). « En général » une conique passe
par cinq points ; on s'attend donc, quand le polygone a au moins cinq sommets
sur chaque fibre (ni triangle ni carré : $\alpha(\alpha^2 - 4)(\alpha + 1)$
inversible), à ce qu'une seule conique passe par tous ses sommets. Pour $Q =
ax^2 + 2bxy + cy^2 + ux + vy + \lambda$ et les cinq points $(\zeta^i, \zeta^{-i})$,
$-2 \leq i \leq 2$, on obtient un système de $5$ équations à $6$ inconnues dont
deux colonnes coïncident (celles de $xy$ et de la constante) ; il est de rang $5$
si et seulement si le déterminant $5 \times 5$ obtenu en supprimant l'une d'elles,
$\Delta(\zeta) \in \mathbf{Z}[\zeta, \zeta^{-1}]$, est inversible. La page voit
qu'il s'annule pour $\zeta^3 = 1$ et $\zeta^4 = 1$, devine la divisibilité par
$\alpha(\alpha^2 - 4)(\alpha + 1)$ et demande : « Étudier la multiplicité dudit
déterminant ? » Le calcul répond :
\[
  \Delta(\zeta) = \pm\, (\alpha - 2)^5\, (\alpha + 2)^2\, \alpha\, (\alpha + 1)^2 ,
\]
de sorte que $\Delta(\zeta)$ est inversible exactement quand $\alpha(\alpha^2 -
4)(\alpha + 1)$ l'est\footnote{Calcul de l'édition, fait par machine ; le signe
dépend de l'ordre des colonnes. La page écrit successivement
$(\zeta^4 - 1)(\zeta^3 - 1)/(\zeta - 1)$ et $(\zeta^2 - 1)(\zeta^3 - 1)/(\zeta - 1)$,
puis identifie $(\zeta^2 + 1)(\zeta^2 - 1)^2(1 + \zeta + \zeta^2)$ à
$\alpha(\alpha^2 - 4)(\alpha + 1)$, ce qui est vrai au facteur $\zeta^{-4}$ près,
une unité. La marge suggère une « manœuvre » de Vandermonde sur les colonnes.}.
Alors (page 13) un polynôme de degré $\leq 2$ nul sur cinq sommets consécutifs
est nul sur tous les sommets, et c'est un multiple de la forme invariante
principale.

\pagerange{17}{17}
\subsection*{Les formes sur l'espace des translations (page 17)}

La page 17 achève la première rédaction. Les formes quasi-invariantes sur
l'espace des translations $\check{V}$ (toujours pour $\alpha^2 - 4$ inversible) sont
parmi celles sur $\check{E}$ : si $\alpha$ est inversible, ce sont les multiples
de $xy$, non dégénérée, qui donne la conique circonscrite ; si $\alpha = 0$
(donc $2$ inversible), ce sont les multiples de $x^2 + y^2$ et de $x^2 - y^2$ ; la
première prend sur les sommets les deux valeurs distinctes $2$ et $\zeta^2 +
\zeta^{-2} = \alpha^2 - 2 = -2$, la seconde y est nulle mais définit une conique
dégénérée. « Il reste à étudier » les cas $\alpha = \pm 2$ et le cas d'un $\alpha$
« variable ».

\pagerange{18}{18}
\subsection*{Seconde rédaction : coordonnées universelles (page 18)}

« Coniques circonscrites à un polygone régulier. » On place trois sommets
consécutifs en $s_{-1} = (0, 1)$, $s_0 = (0, 0)$, $s_1 = (1, 0)$ ; les suivants
sont alors $s_2 = (1 + \alpha, 1)$ et $s_{-2} = (1, 1 + \alpha)$ — les coordonnées
universelles du dossier 78. Pour $q = ax^2 + bxy + cy^2 + ux + vy + d$, les
conditions $q(s_0) = q(s_1) = q(s_{-1}) = 0$ donnent $d = 0$, $u = -a$, $v = -c$,
donc $q = a(x^2 - x) + c(y^2 - y) + bxy$\footnote{La page conclut « $v = -b$ »,
sur un coefficient surchargé ; la forme qu'elle en tire demande $v = -c$.} ; puis
\[
  q(s_2) = (1 + \alpha)(\alpha a + b), \qquad q(s_{-2}) = (1 + \alpha)(\alpha c + b).
\]
Si $1 + \alpha$ est inversible (pas de triangle) il vient $b = -\alpha a =
-\alpha c$, et si de plus $\alpha$ est inversible (pas de carré), $a = c$ : $q$ est
un multiple de
\[
  (2) \qquad q_0(x, y) = (x^2 - x) + (y^2 - y) - \alpha xy ,
\]
« la conique universelle »\footnote{La page écrit « si $\alpha - 1$ inv.\ (i.e.\
on n'a de triangle sur aucune fibre) » ; le facteur en jeu est $1 + \alpha$, et le
triangle est $\alpha = -1$. Une note en coin, peu lisible, propose une
vérification de $q(s_2)$ et $q(s_{-2})$.}. Elle est invariante sous le groupe
diédral.

\pagerange{13}{16}
\subsection*{La forme principale et ses directions isotropes (pages 13 à 16)}

La page 13 poursuit la seconde rédaction\footnote{Ses numéros d'équations ne
prolongent pas ceux des pages 7 à 12, et sa forme $q_0$ est celle de la page
18 ; la page 14 repart de (7) à la suite de la page 13, et la page 19 prolonge
les pages 14 et 15. L'ordre de lecture est donc 18, 13, 14, 15, 16, 19.}. Sur
l'espace des translations, la forme principale est $q_0(x, y) = x^2 + y^2 -
\alpha xy$, de discriminant $\alpha^2 - 4$, inversible exactement dans le cas I.
Homogénéisée,
\[
  q_0(x, y, z) = x^2 + y^2 - \alpha xy - xz - yz, \qquad
  \text{de matrice} \quad
  \begin{pmatrix} 2 & -\alpha & -1 \\ -\alpha & 2 & -1 \\ -1 & -1 & 0 \end{pmatrix},
\]
elle a pour déterminant $-2(\alpha + 2)$ et pour discriminant divisé $-(\alpha +
2)$ : la conique circonscrite est lisse sur chaque fibre où $\alpha \neq -2$,
c'est-à-dire hors du cas III.

Les deux points isotropes à l'infini, confondus exactement quand $\alpha = \pm 2$,
sont les directions de pente $\zeta$ avec $\zeta^2 - \alpha\zeta + 1 = 0$. Se donner
l'une d'elles, c'est remplacer $\alpha$ par la racine inversible $\zeta$, un
paramètre « modulaire » plus fin ; pour les $n$-gones on aura $\Phi_n(\zeta) = 0$
et des conditions de caractéristique. Alors
\[
  q_0(x, y) = (x - \zeta y)(x - \zeta^{-1}y), \qquad
  q(x, y, z) = (x - \zeta y)(x - \zeta^{-1}y) - (x + y)z ,
\]
et l'on veut identifier la conique $C : q = 0$ à $\mathbb{P}(E/D_\zeta)$ en projetant
depuis son point $d_\zeta$, où $D_\zeta = \mathcal{O}_S(1, \zeta, 0)$. La page 15
choisit à cet effet le repère $f_1 = (1, \zeta, 0)$, $f_2 = (1, 0, 1 - \zeta)$, $f_3
= (0, 0, 1)$, et inverse les formules. Une tentative de traduire la donnée
« polygone épinglé par $D_\zeta$ » en termes de fibrés affines et de construire
le schéma en groupes des automorphismes de $P$ qui fixent les deux points à
l'infini est biffée (pages 15 et 16)\footnote{Le passage biffé affirme que ce
schéma en groupes $G$ est lisse de dimension relative $1$, de composante neutre
commutative ; la page 16 est blanche au-delà.}.

\pagerange{19}{19}
\subsection*{Le plan tangent le long de $D_\zeta$ (page 19)}

Le plan tangent à $C$ le long de $D_\zeta$ s'obtient en appliquant la matrice de
$q$ au vecteur $(1, \zeta, 0)$ :
\[
  A = 1 - \zeta^2, \qquad B = \zeta - \zeta^{-1} = -\zeta^{-1}(1 - \zeta)(1 + \zeta),
  \qquad C = -(1 + \zeta)
\]
\footnote{La page écrit d'abord $B = \zeta^{-1}(1 - \zeta^2)$, de signe opposé,
puis la bonne factorisation ; l'équation qu'elle en tire est juste.}. Comme $1
+ \zeta$ est inversible ($\alpha \neq -2$), on divise :
\[
  H_\zeta : \quad z = (1 - \zeta)(x - \zeta^{-1}y), \qquad D_\zeta \subset H_\zeta
  \subset E .
\]
La page en vient alors à identifier $C$ à la droite projective de la droite
affine extension de $E/H_\zeta$ par $H_\zeta/D_\zeta$, scindée : la conique privée de
$d_\zeta$ est une droite affine. La phrase, très surchargée, s'interrompt, et la
liasse avec elle.

\section*{III. Extensions quadratiques (pages 20 à 63)}

La couverture de la page 20 porte « Extensions quadratiques » et « (43~p) ». Les
pages 21 à 63 forment une rédaction suivie, numérotée en paragraphes 1 à 10, avec
une numérotation continue des formules de (1) à (114) ; le § 10 est récrit deux
fois.

\pagerange{21}{22}
\subsection*{1. L'enveloppe affine d'un revêtement quadratique (pages 21 et 22)}

Un \emph{revêtement quadratique} de $S$ est un $S$-schéma $\Omega$ fini et localement
libre de rang $2$ ; c'est la même chose qu'une \emph{algèbre quadratique} $B$
(commutative, unitaire, localement libre de rang $2$), par $\Omega =
\mathrm{Spec}(B)$. Le quotient $B/\mathcal{O}_S$ est inversible ; on note $L$ son dual,
de sorte que $B$ est une extension
\[
  (2) \qquad 0 \to \mathcal{O}_S \xrightarrow{\ i\ } B \to L^{\vee} \to 0
\]
et, dualement,
\[
  (3) \qquad 0 \to L \to B^{\vee} \xrightarrow{\ \pi\ } \mathcal{O}_S \to 0, \qquad
  \pi(\ell) = \ell(e).
\]
Le fibré affine $E_\Omega = \pi^{-1}(1)$ des formes linéaires $\ell : B \to
\mathcal{O}_S$ telles que $\ell(e) = 1$ est un torseur sous $L$ ; c'est $E_\Omega =
\mathrm{Spec}\bigl(\mathrm{Sym}(B)/(e - 1)\bigr)$, et ses sections sont les
scindages de (3). Parmi elles se trouvent les homomorphismes d'algèbres $B \to
\mathcal{O}_S$, d'où une immersion fermée canonique, compatible aux changements de
base,
\[
  (6) \qquad \Omega \hookrightarrow E_\Omega .
\]
Elle est universelle : si $F$ est un fibré affine quelconque, défini par une
extension $0 \to V \to C^{\vee} \to \mathcal{O}_S \to 0$, de sorte que $F =
\mathrm{Spec}(\mathrm{Sym}(C)/(e_C - 1))$, les morphismes $\Omega \to F$ sont les
applications linéaires $C \to B$ envoyant $e_C$ sur $e$, c'est-à-dire, par
transposition, les morphismes de fibrés affines $E_\Omega \to F$ :
\[
  (9) \qquad \mathrm{Hom}_S(\Omega, F) \simeq \mathrm{Homaff}_S(E_\Omega, F).
\]
D'où le nom d'\emph{enveloppe affine} de $\Omega$.

\pagerange{23}{25}
\subsection*{2--3. Scindages, trace, norme et discriminant (pages 23 à 25)}

Une section $s$ de $E_\Omega$, c'est-à-dire un scindage $B \simeq \mathcal{O}_S \oplus
L^{\vee}$, ramène la structure d'algèbre à la multiplication $L^{\vee} \otimes
L^{\vee} \to \mathcal{O}_S \oplus L^{\vee}$, donnée par deux sections $b \in \Gamma(L)$,
$c \in \Gamma(L^{\otimes 2})$ :
\[
  (13) \qquad s(u)\, s(v) = -c\,uv \oplus s(b\,uv), \qquad u, v \in L^{\vee}
\]
\footnote{Une note en marge s'inquiète de la convention de signe. Elle est
cohérente avec (24) ci-dessous.}. La multiplication par $s(u)$ a pour matrice
$\begin{pmatrix} 0 & -cu \\ u & bu \end{pmatrix}$, d'où
\[
  (14)\text{--}(15) \qquad \mathrm{tr}(\lambda + s(u)) = 2\lambda + bu, \qquad
  N(\lambda + s(u)) = \lambda^2 + \lambda bu + cu^2
\]
\footnote{La page porte $-bu$ dans la première matrice et $+bu$ dans la seconde ;
les traces qu'elle en tire demandent $+bu$.}. Dans $\check{\mathbb{V}}(L) \simeq
E_\Omega$, le revêtement est le sous-schéma d'équation $x^2 - bx + c = 0$, et la
catégorie des $\Omega$ munis d'un scindage équivaut à celle des triples $(L, b,
c)$ (notes en marge).

Changer de scindage, $s' = s + \tau$ avec $\tau \in \Gamma(L)$, change $(b, c)$ en
\[
  (17) \qquad b' = b + 2\tau, \qquad c' = c + \tau b + \tau^2 .
\]
Le \emph{discriminant}
\[
  (18) \qquad \mathfrak{D} = \mathfrak{D}(\Omega/S) = b^2 - 4c \in \Gamma(S,
  L^{\otimes 2})
\]
est donc indépendant du scindage. La forme polaire de la norme est
\[
  (20) \qquad \varphi(\xi, \eta) = \mathrm{tr}\,\xi\;\mathrm{tr}\,\eta -
  \mathrm{tr}(\xi\eta),
\]
d'où $\varphi(\xi, e) = \mathrm{tr}\,\xi$, et, en calculant $\Lambda^2\varphi(e \wedge
\xi, e \wedge \xi) = (\mathrm{tr}\,\xi)^2 - 2\,\mathrm{tr}(\xi^2)$ pour $\xi = s(u)$,
\[
  (19) \qquad \delta(N_{\Omega/S}) = -\mathfrak{D}(\Omega/S),
\]
le déterminant de la forme polaire de la norme étant lu comme section de
$L^{\otimes 2}$ par $\det B \simeq L^{\vee}$\footnote{La page a d'abord écrit
$-4\mathfrak{D}$, puis biffé le $4$ ; c'est la forme biffée qui serait fausse. Une
première tentative, encadrée et barrée, figure au bas de la page 24.}.

\pagerange{26}{26}
\subsection*{4. Une base $(1, t)$ (page 26)}

Si $L$ est trivialisé et le scindage donné, c'est-à-dire si $E_\Omega$ est
identifié à la droite affine standard, on dispose d'une section $t$ de $B$ telle
que $(1, t)$ soit une base ; alors
\[
  (24) \qquad t^2 - bt + c = 0, \qquad (25) \qquad b = \mathrm{tr}\,t, \quad c =
  N(t),
\]
et $\mathfrak{D} = b^2 - 4c$\footnote{(24) se déduit de (13) avec $u = v$ égal au
générateur : $t^2 = -c + bt$. La note de la transcription selon laquelle les
signes de $c$ diffèrent entre (13) et (24) ne se confirme pas : ils
concordent.}.

\pagerange{27}{33}
\subsection*{5. Revêtements quadratiques et puissances symétriques d'une droite affine (pages 27 à 33)}

Un revêtement quadratique est la même chose qu'un fibré affine de rang $1$ muni
d'un diviseur relatif de degré $2$ : $\Omega \subset E_\Omega$ en est un, et
réciproquement un diviseur relatif de degré $2$ d'une droite affine $F$ a $F$ pour
enveloppe affine. Or les diviseurs relatifs effectifs de degré $n$ d'une droite
affine $F$ sont les sections de la puissance symétrique $\Sigma^n(F) =
F^n/\mathfrak{S}_n$, et celle-ci est un fibré affine de rang $n$. Quand $F$ est
trivialisé par une section $s$, $F \simeq \check{\mathbb{V}}(L)$, les fonctions
symétriques élémentaires donnent
\[
  (28) \qquad \Sigma^n\check{\mathbb{V}}(L) \xrightarrow{\ \sim\ }
  \check{\mathbb{V}}\Bigl(\bigoplus_{1 \leq i \leq n} L^{\otimes i}\Bigr), \qquad
  (x_1, \ldots, x_n) \mapsto (\sigma_i(x))_{1 \leq i \leq n}.
\]
La translation par $\tau \in \Gamma(L)$ y devient la transformation affine
\[
  (33) \qquad \varphi_n(\tau)_i(\sigma_1, \ldots, \sigma_n) = \sum_{p + q = i}
  \binom{n - p}{q} \tau^q \sigma_p \qquad (\sigma_0 = 1),
\]
d'où, en recollant, une structure de fibré affine sur $\Sigma^n(F)$ (34). Pour
$n = 2$,
\[
  (35) \qquad \varphi_2(\tau)(\sigma_1, \sigma_2) = (\sigma_1 + 2\tau,\ \sigma_2 +
  \tau\sigma_1 + \tau^2),
\]
de partie linéaire $(\sigma_1, \sigma_2) \mapsto (\sigma_1, \sigma_2 + \tau\sigma_1)$.

\textbf{Proposition} (page 30). Si $F$ est le fibré affine de rang $1$ défini par
l'extension (3), $\Sigma^2(F)$ est canoniquement un fibré affine sous le fibré
vectoriel $B^{\vee} \otimes L$, extension de $L$ par $L^{\otimes 2}$, et
\[
  (39) \qquad \Sigma^2(F)/\check{\mathbb{V}}(L^{\otimes 2}) \simeq 2F,
\]
où $2F$ est le $L$-torseur déduit de $F$ par la multiplication par $2$ (dans
$\check{\mathbb{V}}(B^{\vee})$, l'image réciproque de la section $2$). Le morphisme
\[
  (40) \qquad \underline{b} : \Sigma^2(F) \to 2F
\]
est caractérisé par sa compatibilité aux changements de base et par
$\underline{b}(\Omega(s_1, s_2)) = s_1 + s_2$, où $\Omega(s_1, s_2)$ est le diviseur
défini par deux sections : c'est une « trace universelle ». Pour toute section
$s$ de $F$, le $b$ du scindage $s$ est
\[
  (42) \qquad b(s) = \underline{b} - 2s ,
\]
ce qui donne bien $(s_1 - s) + (s_2 - s)$ pour $\Omega(s_1, s_2)$.

Se donner une section de $\Sigma^2(F)$, c'est donc se donner successivement une
section $\underline{b}$ de $2F$, puis une section du $L^{\otimes 2}$-torseur
$\Sigma(\underline{b})$, fibre de $\Sigma^2(F) \to 2F$. D'où deux obstructions :
$2F$ a une section si et seulement si $2\,\mathrm{cl}(F) = 0$ dans $H^1(S, L)$, et ses
sections forment alors un torseur sous $H^0(S, L)$ ; $\underline{b}$ choisi, les
relèvements existent si et seulement si $\mathrm{cl}(\Sigma(\underline{b})) = 0$
dans $H^1(S, L^{\otimes 2})$, et forment un torseur sous $H^0(S, L^{\otimes 2})$.
Enfin (45)--(46) : $\Sigma^2(F)$ est le fibré affine associé à l'extension
$\mathrm{Sym}^2(B^{\vee})$ de $\mathcal{O}_S$ par $B^{\vee} \otimes L$\footnote{La page
écrit, dans la justification et dans (46), « extension de $\mathcal{O}_S$ par
$V^{(1)} = B^{\vee}$ » ; le noyau de $\mathrm{Sym}^2(B^{\vee}) \to
\mathrm{Sym}^2(\mathcal{O}_S)$ est l'image de $L \otimes B^{\vee}$, comme le dit
d'ailleurs (45).}.

\pagerange{34}{40}
\subsection*{5 bis. Seconde mouture : $\Sigma^n(F)$ en général, et le théorème (pages 34 à 40)}

« Détermination de $\Sigma^n(F)$ (2ème mouture). » Soit $F = \pi_1^{-1}(1)$ défini
par $0 \to L \to V \xrightarrow{\pi_1} \mathcal{O}_S \to 0$, avec $V = B^{\vee}$.
Pour tout $n \geq 1$, on a une extension
\[
  (49) \qquad 0 \to L \otimes \mathrm{Sym}^{n-1}(V) \xrightarrow{\ i_n\ }
  \mathrm{Sym}^n(V) \xrightarrow{\ \pi_n\ } \mathcal{O}_S \to 0,
\]
et une filtration décroissante de $\mathrm{Sym}^n(V)$ par les $L^{\otimes i} \otimes
\mathrm{Sym}^{n-i}(V)$ (50). Comme $V$ est de rang $2$, $V^{\vee} \simeq V \otimes
(\det V)^{-1}$, et l'extension fournit $\kappa : \det V \xrightarrow{\sim} L$,
$\kappa(i_1(\lambda) \wedge x) = \lambda\,\pi_1(x)$ (55) ; d'où $B = V^{\vee} \simeq V
\otimes L^{-1}$, l'extension (59) $0 \to \mathcal{O}_S \to L^{-1} \otimes V \to
L^{\vee} \to 0$ s'identifiant à (2)\footnote{La page appelle $\alpha$
l'isomorphisme $\kappa$ ; on change la lettre, réservée à $\zeta + \zeta^{-1}$. Un
premier départ de la page 34, sur $\Sigma^n F = \pi_n^{-1}(1)$, est biffé.}.

\textbf{Proposition 61.} Pour tout $n$, l'application $J_n :
\mathrm{Sym}^n(B) \to \mathrm{Algaff}(F) = \mathrm{Sym}(B)/(e - 1)$ est injective, son
image est formée des fonctions polynômes de degré $\leq n$, les $J_n$ sont
compatibles à la multiplication par $e$, et $\mathrm{Algaff}(F)$ est la limite
inductive des $\mathrm{Sym}^n(B) \simeq L^{\otimes -n} \otimes \mathrm{Sym}^n(V)$. Le
quotient $\mathrm{Sym}^n(B)/\mathrm{Sym}^{n-1}(B) \simeq L^{\otimes -n}$ reçoit le
\emph{terme dominant} d'une fonction de degré $\leq n$.

\textbf{Proposition} (page 38). Pour un diviseur relatif $D$ de degré $n$ sur $F$,
le sous-Module $M_D \subset \mathrm{Sym}^n(B)$ des fonctions de degré $\leq n$ nulles
sur $D$ est un supplémentaire de $\mathrm{Sym}^{n-1}(B)$ ; $D \mapsto M_D$ est une
bijection, compatible aux changements de base, sur les tels supplémentaires,
d'où
\[
  \underline{\mathrm{Div}}^{+n}_{F/S} \simeq \Sigma^n(F/S) \simeq
  \underline{\mathrm{Scind}}\bigl(0 \to \mathrm{Sym}^{n-1}(B) \to \mathrm{Sym}^n(B) \to
  L^{\otimes -n} \to 0\bigr).
\]
($M_D$ est engendré par l'équation « unitaire » de $D$.) Par dualité,
\[
  (62) \qquad \Sigma^n(F/S) \simeq \text{fibré affine associé à } 0 \to
  \mathrm{Sym}^{n-1}(V) \otimes L \to \mathrm{Sym}^n(V) \to \mathcal{O}_S \to 0,
\]
la flèche étant la valeur en $e$ : \emph{un diviseur relatif de degré $n$ sur $F$
est la même chose qu'une forme de degré $n$ sur $B$ valant $1$ en $e$.}

\textbf{Théorème} (§ 6, pages 39 et 40). Soit $B$ un Module localement libre de
rang $2$ muni d'une section $e$ partout non nulle, et $F$ la droite affine
correspondante. Il y a des bijections, compatibles aux changements de base,
entre
\begin{itemize}
  \item[a)] les structures d'algèbre commutative unitaire sur $B$ d'unité $e$ ;
  \item[b)] les diviseurs relatifs de degré $2$ de $F$ ;
  \item[c)] les sections de $\Sigma^2(F/S)$ ;
  \item[d)] les formes quadratiques $q$ sur $B$ telles que $q(e) = 1$,
\end{itemize}
et la forme $q$ attachée à une structure d'algèbre est sa norme : $q = N_{B/\mathcal{O}_S}$
(65). (Directement : $q$ étant donnée, $\mathrm{tr}(x) = \varphi(e, x)$ et $x^2 =
\mathrm{tr}(x)\,x - q(x)\,e$ définissent l'algèbre.)

Pour vérifier (65), la page passe au rang $n$ quelconque : une $n$-forme $f_n$
sur $B$ avec $f_n(e) = 1$ équivaut à une algèbre $C_n$ localement libre de rang
$n$ munie d'une application linéaire $g : B \to C_n$ avec $g(e) = 1$ et $g(B)$
engendrant $C_n$ (66)--(67)\footnote{La page écrit « a) $q(e_B) = 1_{C_n}$ » ; il
s'agit de $g(e_B) = 1_{C_n}$.}, et l'on a « sauf erreur »
\[
  (70) \qquad f_n = N_{C_n/\mathcal{O}_S} \circ g .
\]
Si $(e, t)$ est une base de $B$, $C_n = \mathcal{O}_S[T]/(f)$ avec $f = T^n + a_1T^{n-1}
+ \cdots + a_n$ unitaire, et
\[
  (74) \qquad N_{C_n/\mathcal{O}_S}(u + vt) = u^n - a_1u^{n-1}v + \cdots +
  (-1)^n a_n v^n = (-v)^n f(-u/v),
\]
d'où $N(t) = (-1)^n a_n$, le déterminant de la matrice compagnon. Par un argument
universel, ou par extension fidèlement plate, on se ramène à $f = \prod (T -
t_i)$, où (74) devient $\prod_i (u + vt_i) = \prod_i \rho_i(u + vt)$, cas
particulier de la formule « bien connue » $N(w) = \prod_i \rho_i(w)$ pour
l'algèbre déployée\footnote{Une ligne de la page 42 écrit $(-v)^n f(1 -
\frac{u}{v})$, le « $1$ » étant surchargé, sans doute biffé. Les numéros (73) à
(76) ne sont pas dans l'ordre de la page.}.

\pagerange{41}{44}
\subsection*{La trace, et les formes de trace donnée (pages 41 à 44)}

\textbf{Proposition} (page 43). Pour $q$ comme en d), la forme $\underline{b}(q) :
x \mapsto \varphi(e, x) = \mathrm{tr}_{B/\mathcal{O}_S}(x)$ vérifie $\underline{b}(q)(e) =
2q(e) = 2$ ; c'est une section de $2F$, et c'est le $\underline{b}(\Omega)$ du § 5.

Les formes $q$ de trace donnée forment un torseur sous le Module des formes
$q_0$ telles que $q_0(e) = 0$ et $\varphi_0(e, \cdot) = 0$, c'est-à-dire des formes
qui se factorisent par $B \to L^{\vee}$ : $q_0(x) = c_0\,\pi_1(x)^2$ avec $c_0 \in
\Gamma(L^{\otimes 2})$ (83)\footnote{La transcription lit l'exposant de $L$
« $\otimes{-2}$ », en signalant qu'il pourrait être $\otimes 2$ : une forme
quadratique sur $L^{\vee}$ est une section de $L^{\otimes 2}$, et c'est ce que
demande la suite (pages 55 et 56).}.

\pagerange{45}{47}
\subsection*{7. L'involution canonique (pages 45 à 47)}

\textbf{Proposition.} Pour une algèbre quadratique $B$, l'application
\[
  (84) \qquad \sigma : u \mapsto (\mathrm{tr}\,u)\,e - u
\]
est un automorphisme d'algèbre involutif, caractérisé par $u + \sigma u =
(\mathrm{tr}\,u)\,e$. On a
\[
  (86)\text{--}(87) \qquad u\,\sigma(u) = N(u)\,e, \qquad u^2 - (\mathrm{tr}\,u)\,u +
  N(u) = 0,
\]
l'équation de Hamilton--Cayley de la multiplication par $u$.

La donnée de $\sigma$ équivaut à celle de la trace. Plus précisément,
$-\sigma = \mathrm{id} - \mathrm{tr} \otimes e$ est une pseudo-réflexion de $B$ : l'identité
modulo la droite $\mathcal{O}_S e$, qu'elle envoie sur son opposée ; les
pseudo-réflexions $\mathrm{id} + a \otimes e$ de ce type sont involutives exactement
quand $a(e) = -2$, et $a = -\mathrm{tr}$ redonne $-\sigma$\footnote{La page calcule
$\sigma^2 = \mathrm{id} + (2 + \langle a, \pi_1\rangle)\,\pi_1 \otimes a$ et parle de
$\mathcal{O}_S e$ comme de la « droite fixe » : c'est la droite fixe de $\sigma$, et
celle où la pseudo-réflexion $-\sigma$ agit par $-1$. La marge, écrite le long du
bord, est en grande partie illisible.}. \textbf{Proposition} (page 46) : une
telle pseudo-réflexion est l'identité sur une fibre si et seulement si cette
fibre est de caractéristique $2$ et que la trace y est nulle. Ainsi $\sigma \neq
\mathrm{id}$ sur la fibre en $s$ si et seulement si $k(s)$ n'est pas de
caractéristique $2$, ou $\mathrm{tr}_s \neq 0$.

\pagerange{47}{49}
\subsection*{8. Quand $2$ est inversible (pages 47 à 49)}

Si $2$ est inversible sur $S$, la donnée de $\underline{b} = \mathrm{tr} \in \Gamma(2F)$
équivaut à celle de l'origine $s_0 = \tfrac{1}{2}\underline{b}$ de $F$, c'est-à-dire à
un scindage $B = \mathcal{O}_S \oplus \ker(\mathrm{tr})$, où $\ker(\mathrm{tr})$ est
l'espace propre de $\sigma$ pour $-1$. Pour ce scindage $b(s_0) = 0$ (90), de sorte
que l'algèbre est décrite par $c_0 \in \Gamma(L^{\otimes 2})$ seul, et
$\mathfrak{D} = -4c_0$ (91).

\textbf{Théorème} (page 48). Si $2$ est inversible sur $S$, la catégorie des
algèbres quadratiques sur $\mathcal{O}_S$ (avec leurs isomorphismes) équivaut à celle
des couples $(L, \mathfrak{D})$, $L$ inversible et $\mathfrak{D} \in \Gamma(L^{\otimes
2})$. Le fibré affine de $\Omega$ est canoniquement scindé, $\Omega \subset
\check{\mathbb{V}}(L)$, et si $\Omega(\mathfrak{D}) \subset \check{\mathbb{V}}(L)$ est le
sous-schéma $y^2 = \mathfrak{D}$, on a
\[
  (92) \qquad \Omega = \tfrac{1}{2}\,\Omega(\mathfrak{D}) :
\]
dans le scindage canonique l'équation $x^2 - bx + c = 0$ devient $x^2 =
\tfrac{1}{4}\mathfrak{D}$, soit $(2x)^2 = \mathfrak{D}$.

\textbf{Corollaire} (page 49). $\mu_2$ opère sur $\Omega(\mathfrak{D})$ par $y \mapsto
-y$, et l'opération induite sur $\Omega$ est l'involution canonique du § 7, qui
s'écrit en général $\sigma x = b - x$ (94).

Comme quotient de $\mathrm{Sym}(L^{\vee}) = \bigoplus_n L^{\otimes -n}$, $B$ s'écrit
ici $\mathrm{Sym}(L^{\vee})/\bigl(w - \tfrac{1}{4}\mathfrak{D}(w) ;\ w \in L^{\otimes
-2}\bigr)$\footnote{La page écrit « $B \simeq \mathrm{Sym}(L^{\vee})/(1 -
4\partial)\,L^{\otimes -2}$ », après une formule générale (93) en partie noircie ;
la relation qu'impose (13) avec $b = 0$ et $c = -\mathfrak{D}/4$ est $s(u)^2 =
\tfrac{1}{4}\mathfrak{D}(u^2)$.}.

\pagerange{50}{52}
\subsection*{9. En caractéristique $2$ (pages 50 à 52)}

Si $2 = 0$ sur $S$, alors $2F \simeq \check{\mathbb{V}}(L)$ canoniquement, la condition
$\mathrm{tr}(e) = 2$ devient $\mathrm{tr}(e) = 0$, et la trace est une section
\[
  (97) \qquad \underline{b} = \mathrm{tr} \in \Gamma(S, L),
\]
égale au $b(s)$ de tout scindage (99) ; le discriminant est
\[
  (98) \qquad \mathfrak{D} = \underline{b}^2 .
\]
L'équation $x^2 = \mathfrak{D}$ s'écrit $(x - \underline{b})^2 = 0$ et n'a plus « grand
rapport » avec $\Omega$. La classe $\gamma \in H^1(S, L)$ du torseur $F$ est un
invariant important de la situation (100). L'involution canonique est la
translation par $\underline{b}$ ; c'est une réflexion (non seulement une
pseudo-réflexion) si et seulement si $\underline{b}$ est inversible, si et seulement
si $\mathfrak{D}$ l'est, si et seulement si $\Omega$ est étale.

Dans ce cas, pour $\underline{b}$ fixé, on a la suite exacte d'Artin--Schreier
\[
  (101) \qquad 0 \to (\mathbf{Z}/2)_S \to \check{\mathbb{V}}(L) \xrightarrow{\ \beta\ }
  \check{\mathbb{V}}(L^{\otimes 2}) \to 0, \qquad \varepsilon \mapsto
  \varepsilon\underline{b}, \quad x \mapsto x^2 + \underline{b}\,x,
\]
et le torseur $C = F \wedge^{\check{\mathbb{V}}(L), \beta} \check{\mathbb{V}}(L^{\otimes
2})$ sous $\check{\mathbb{V}}(L^{\otimes 2})$, dont $F$ est un revêtement étale
quadratique. Les $\Omega$ de trace $\underline{b}$ correspondent biunivoquement aux
sections $\sigma$ de $C$ : $\Omega$ est l'image réciproque par $\sigma$ du revêtement
$F \to C$ — dans un scindage, l'équation $x^2 + \underline{b}x = c$\footnote{La
page 52 note qu'« il faudrait préciser » que la structure de torseur sous
$L^{\otimes 2}$ sur les $q$ compatibles à $\underline{b}$ est celle des $\sigma$. Le
dossier 84 (pages 56 et 57) énonce le même théorème en caractéristique $2$, avec
les données $(L, P_1, T, \alpha)$.}.

\pagerange{52}{57}
\subsection*{10. Le cas mixte, première rédaction (pages 52 à 57)}

On suppose désormais que $2$ est \emph{régulier} sur $S$ (non diviseur de zéro
dans $\mathcal{O}_S$), et l'on pose $S_0 = V(2\mathcal{O}_S)$, $S_1 = V(4\mathcal{O}_S)$, son
premier voisinage infinitésimal. Le point de départ est le cas $S = S_2 \cup
S_{\mathfrak{D}}$ — $2$ inversible ou $\Omega$ étale en chaque point —, où l'on doit
recoller la description du § 8 sur $S_2$ et un revêtement étale sur
$S_{\mathfrak{D}}$ ; la page passe vite au cas général.

\textbf{Le torseur $F$ et la section $b_0$.} Un $L$-torseur $F$ muni d'une section
$\underline{b}$ de $2F$ est trivialisé, au-dessus de $S[1/2]$, par
$\tfrac{1}{2}\underline{b}$. Pour toute section locale $s$ de $F$, $b(s) =
\underline{b} - 2s \in L$ (42) change de $2\tau$ quand $s$ change de $\tau$ ; sa
réduction
\[
  (105) \qquad b_0 = b(s)|S_0 \in \Gamma(S_0, L_0)
\]
est donc globale, et c'est la trace $\underline{b}(\Omega_0/S_0)$ du § 9.
Réciproquement, $2$ étant régulier, $(F, \underline{b})$ se reconstruit à partir de
$(L, b_0)$ : $F$ est le faisceau des sections $y$ de $\tfrac{1}{2}L$ telles que $2y$
soit une section de $L$ relevant $b_0$\footnote{La page le dit en termes de
$\underline{L} \subset \tfrac{1}{f}\underline{L} \subset i_*(L|S_f)$, $f = 2$, et de la
section $x_0$ de $\tfrac{1}{f}L/L \simeq L/fL$ ; plusieurs mots de ces lignes sont
incertains, et une note marginale réserve le cas d'un $f$ seulement
$L$-régulier.}.

La condition
\[
  (106) \qquad b_0^2 = \mathfrak{D}_0 = \mathfrak{D}|S_0
\]
est évidemment nécessaire ; si $S_0$ est réduit, elle détermine $b_0$, une racine
carrée dans un anneau réduit de caractéristique $2$ étant unique. Les formes $q$ de
trace $\underline{b}$ forment un torseur $Q(\underline{b})$ sous
$\check{\mathbb{V}}(L^{\otimes 2})$ (page 44), et l'application $\Theta : q \mapsto
\mathfrak{D}(q)$ est affine :
\[
  (107)\text{--}(108) \qquad \Theta(q + \varpi) = \Theta(q) - 4\varpi, \qquad
  \Theta_0 = -4\,\mathrm{id}.
\]
La forme cherchée vérifie $\Theta(q) = \mathfrak{D}$ ; pour $q = q_0 + \varpi$, cela
s'écrit $4\varpi = \Theta(q_0) - \mathfrak{D}$ (111). Comme $4$ est régulier,
$\varpi$ est unique s'il existe, et il existe si et seulement si la section
\[
  (113) \qquad a_1(\mathfrak{D}) = \bigl(\Theta(q_0) - \mathfrak{D}\bigr)|S_1 \in
  \Gamma(S_1, L_1^{\otimes 2})
\]
est nulle (114). Elle ne dépend pas de $q_0$, puisque changer $q_0$ en $q_0 +
\varpi$ la change de $4\varpi|S_1 = 0$\footnote{La phrase de la page sur la
dépendance en $q_0$ est mal lisible (« il dépend du choix de $q_0$, chaque section
[définie modulo …] ») ; le calcul montre qu'il n'en dépend pas.}.

\textbf{Théorème} (page 56). Soit $S$ un schéma où $2$ est régulier. La catégorie des
revêtements quadratiques de $S$ équivaut à celle des triples $(L, \mathfrak{D},
b_0)$, $L$ inversible, $\mathfrak{D} \in \Gamma(L^{\otimes 2})$, $b_0 \in \Gamma(S_0,
L_0)$, satisfaisant
\begin{itemize}
  \item[a)] $b_0^2 = \mathfrak{D}_0$ ;
  \item[b)] $a_1(\mathfrak{D}) = 0$.
\end{itemize}
Quand $S_0$ est réduit, $b_0$ est déterminé par $(L, \mathfrak{D})$, et a) devient une
condition d'existence d'une racine carrée de $\mathfrak{D}_0$\footnote{Comme $\Theta(q_0)
\equiv b^2 \pmod 4$ pour tout relèvement $b$ de $b_0$, on a $a_1(\mathfrak{D}) =
\gamma(b_0) - \mathfrak{D}_1$, où $\gamma(b_0)$ est le carré modulo $4$ d'un relèvement
quelconque de $b_0$ (voir la seconde rédaction) ; b) implique donc a). La page les
énonce séparément.}.

\textbf{Corollaire 1} (page 57). Soit $(L, \mathfrak{D}, b_0)$ satisfaisant a), et $U
\subset S$ un ouvert au-dessus duquel il existe un revêtement $\Omega_U$ correspondant.
Si $U \cap S_0$ est schématiquement dense dans $S_0$ (par exemple $S_0$ réduit et $U
\cap S_0$ dense), il existe $\Omega$ sur $S$ correspondant à $(L, \mathfrak{D},
b_0)$\footnote{La page ajoute entre crochets « et $U \cap S_1$ schématiquement dense
dans $S_1$ ». Cela découle de la densité de $U \cap S_0$ dans $S_0$, $2$ étant
régulier : une section de $L_1^{\otimes 2}$ nulle sur $U$ a une image nulle dans
$L_0^{\otimes 2}$, donc est dans $2L^{\otimes 2}/4L^{\otimes 2} \simeq
L_0^{\otimes 2}$, où elle est encore nulle sur $U$.}. En effet la condition b) est
une égalité de sections de $L_1^{\otimes 2}$ sur $S_1$, vraie sur $U \cap S_1$.

\textbf{Corollaire 2} (page 57). Supposons $S_0$ réduit et normal. Pour qu'il existe un
revêtement quadratique de $S$ correspondant à $(L, \mathfrak{D})$, il faut et il suffit
que le problème ait une solution sur $\mathrm{Spec}\,\mathcal{O}_{S, s}$ pour tout point
maximal $s$ de $S_0$. L'énoncé s'arrête sur des points de suspension ; il est repris
et justifié page 63 (Corollaire 3 ci-dessous).

\pagerange{58}{60}
\subsection*{10. Le cas mixte, seconde rédaction (pages 58 à 60)}

\textbf{Théorème.} Soit $S$ un schéma où $2$ est régulier. À tout revêtement
quadratique $\Omega$ on associe $(L, \mathfrak{D}, b_0)$ : $L = (\det
\mathrm{Aff}(\Omega/S))^{-1}$, $\mathfrak{D} = \mathfrak{D}(\Omega/S)$, $b_0 =
\underline{b}(\Omega_0/S_0)$, de sorte que $b_0^2 = \mathfrak{D}_0$.
\begin{itemize}
  \item[a)] Le foncteur $\Omega \mapsto (L, \mathfrak{D}, b_0)$, entre groupoïdes, est
    pleinement fidèle.
  \item[b)] Si $b_0$ se relève en une section de $L$ sur $S$ — ce qui est le cas si $S$
    est affine, et plus généralement si la multiplication par $2$ est injective sur
    $H^1(S, L)$ —, alors $(L, \mathfrak{D}, b_0)$ est dans l'image essentielle si et
    seulement si $\mathfrak{D} = b^2 - 4c$ avec $b \in \Gamma(S, L)$ relevant $b_0$ et $c
    \in \Gamma(S, L^{\otimes 2})$.
  \item[c)] En tout cas, $(L, \mathfrak{D}, b_0)$ est dans l'image essentielle si et
    seulement si
    \[
      \mathfrak{D}_1 = \gamma(b_0),
    \]
    où $\gamma : L_0 \to L_1^{\otimes 2}$ est l'application « carré d'un relèvement »,
    bien définie puisque $\beta_1 \equiv \beta'_1 \pmod 2 \Rightarrow \beta_1^2 \equiv
    \beta_1'^2 \pmod 4$.
\end{itemize}
\footnote{La condition de b) est sur la page « si $H^1(S, L) \to H^1(S \smallsetminus
S_0, L)$ est injectif (p.\ ex.\ si $S$ affine) », puis, page 60, « $H^1(S, L) \to
H^1(S, L_0)$ injectif ». Le point est que $F$ est trivial si et seulement si $b_0$ se
relève en une section globale de $L$ ; par la suite $0 \to L \xrightarrow{2} L \to L_0
\to 0$, l'obstruction est dans le noyau de la multiplication par $2$ sur $H^1(S, L)$.
La condition de la page sur $S \smallsetminus S_0$, où $2$ est inversible, implique
l'injectivité de cette multiplication ; on énonce la condition elle-même. La page 60
écrit « $(L, \partial, c_0)$ » où le contexte demande $b_0$.}

\emph{Démonstration de a).} On peut supposer $S$ affine ; alors les $\Omega$ sont les
$\Omega(b, c) = \{x \mid x^2 - bx + c = 0\} \subset \check{\mathbb{V}}(L)$, avec
$b|S_0 = b_0$ et $b^2 - 4c = \mathfrak{D}$, et il faut voir qu'entre deux d'entre eux
il y a un isomorphisme et un seul. Les isomorphismes sont les $\tau \in \Gamma(L)$
avec $b' = b + 2\tau$, $c' = c + \tau b + \tau^2$ (17). Comme $b$ et $b'$ ont même
réduction modulo $2$, $b' - b = 2\tau$ pour un unique $\tau$ ($2$ régulier) ; et
$c'_1 = c + \tau b + \tau^2$ vérifie $b'^2 - 4c'_1 = \mathfrak{D} = b'^2 - 4c'$, donc
$4(c'_1 - c') = 0$ et $c' = c'_1$ ($4$ régulier). \emph{b)} : si $F$ est trivial, les
$\Omega$ sont des $\Omega(b, c)$, et ces conditions impliquent $\mathfrak{D}_1 = b_1^2$
avec $b_1 = b|S_1$ ; inversement, si $b$ relève $b_0$ et $b^2 \equiv \mathfrak{D}
\pmod 4$, on écrit $b^2 - \mathfrak{D} = 4c$. \emph{c)} : la condition d'existence
devient locale par a), et localement c'est b), puisque $b^2|S_1 = \gamma(b_0)$ pour
tout relèvement $b$ de $b_0$. Le dossier 84 (pages 48 à 53) énonce le même
théorème sous la forme $(L, \delta, T_0)$, $T_0^2 \equiv \delta \pmod 4$.

\pagerange{61}{63}
\subsection*{Troisième rédaction : les corollaires (pages 61 à 63)}

La page 61 reprend au milieu de la démonstration de c), avec la lettre $\vartheta$
pour le discriminant.

\textbf{Corollaire 1.} $(L, \mathfrak{D}, b_0)$ est dans l'image essentielle s'il
existe $b_1 \in \Gamma(S_1, L_1)$ relevant $b_0$ avec $b_1^2 = \mathfrak{D}_1$, et
seulement si, lorsque $H^1(S_0, L_0) = 0$ (par exemple $S_0$ affine) : l'obstruction à
relever $b_0$ en $b_1$ est dans $H^1(S_0, L_0)$, par $0 \to L_0 \xrightarrow{2} L_1
\to L_0 \to 0$.

\textbf{Corollaire 2.} Supposons $S_0$ réduit. Alors $\Omega \mapsto (L, \mathfrak{D})$
est pleinement fidèle, et son image essentielle est formée des $(L, \mathfrak{D})$ tels
que $\mathfrak{D}_1$ soit localement le carré d'une section de $L_1$ ; il suffit, et il
faut si $H^1(S_0, L_0) = 0$, qu'il existe $b_1 \in \Gamma(S_1, L_1)$ avec $b_1^2 =
\mathfrak{D}_1$ ; et il suffit, et il faut si $F$ est trivial, qu'il existe $b, c$ avec
$b^2 - 4c = \mathfrak{D}$. En effet, $S_0$ étant réduit, $\mathfrak{D}_0$ a au plus une
racine carrée.

\emph{NB.} La condition « $\mathfrak{D}_1 = b_1^2$ localement » se décompose en a)
$\mathfrak{D}_0 = b_0^2$ et b) $\mathfrak{D}_1 = \gamma(b_0)$, qui s'écrit $\alpha(\mathfrak{D}_1)
= 0$ où $\alpha(\mathfrak{D}_1) \in \Gamma(S_0, L_0^{\otimes 2})$ est défini par
$\beta(\alpha(\mathfrak{D}_1)) = \gamma(b_0) - \mathfrak{D}_1$, $\beta : L_0^{\otimes 2}
\xrightarrow{\sim} \mathrm{Ker}(L_1^{\otimes 2} \to L_0^{\otimes 2})$ étant induit par la
multiplication par $2$\footnote{La page note $L_0$ et $L_1$ là où il s'agit de leurs
carrés tensoriels. Elle ajoute, pour a), « la classe $\partial(b_0) \in H^1(S_0,
\mu_2)$ est nulle ! » ; nous ne voyons pas de quelle classe il s'agit, $b_0$ étant ici
une section déjà donnée.}.

\textbf{Corollaire 3} (page 63). Supposons $S_0$ réduit et normal. Pour que $(L,
\mathfrak{D})$ provienne d'un revêtement quadratique de $S$, il faut et il suffit qu'il en
soit ainsi aux points maximaux de $S_0$. En effet, une racine carrée de $\mathfrak{D}_0$
aux points maximaux est entière sur $\mathcal{O}_{S_0}$, donc s'étend en $b_0$ par
normalité ; et $\alpha(\mathfrak{D}_1)$, nul aux points maximaux de $S_0$, est nul
puisque $S_0$ est réduit\footnote{Cette justification est de l'édition ; la page
énonce le corollaire sans démonstration.}.

\section*{IV. Discriminant des formes quadratiques (pages 64 à 87)}

La couverture de la page 64 porte « discriminant des formes quadratiques ».

\pagerange{65}{71}
\subsection*{Le discriminant universel comme somme sur des graphes (pages 65 à 71)}

Soit $I$ un ensemble fini, $AQ_I = \mathbf{Z}[Y_J]$ l'anneau de la forme
quadratique universelle indexée par $I$, les indéterminées $Y_J$ étant indexées par
les parties $J \subset I$ de cardinal $1$ ou $2$ : la forme est $\sum_i Y_i x_i^2 +
\sum_{i < j} Y_{ij} x_i x_j$. Sa matrice a pour coefficients $2Y_i$ sur la diagonale
et $Y_{ij}$ ailleurs ; son déterminant $\Delta$ est le discriminant universel.
Développé selon les permutations, $\Delta = \sum_\sigma \varepsilon_\sigma \prod_i
X_{i, \sigma i}$, chaque $\sigma$ se décompose en points fixes, transpositions et
cycles de longueur $\geq 3$ ; le signe est $\varepsilon_\sigma = (-1)^{\nu_2 +
\nu_*^+}$, où $\nu_2$ est le nombre de transpositions et $\nu_*^+$ celui des cycles de
longueur paire $\geq 4$\footnote{La page définit $I_*(\sigma)$ comme la réunion des
cycles de cardinal $\geq 3$, puis, dans les définitions de $\nu_*^+$ et $\nu_*$, écrit
une borne repassée qui se lit $2$ ; c'est $3$ (et $4$ pour les cycles pairs) qu'il
faut.}.

Le monôme d'une permutation ne dépend que du \emph{graphe} $\gamma$ qu'elle définit
sur $I$ : des sommets isolés (les points fixes), des arêtes doubles (les
transpositions) et des polygones d'ordre $\geq 3$ (les autres cycles, sans leur
orientation). Chaque $\gamma$ provient de $2^{\nu_*(\gamma)}$ permutations, $\nu_*$
étant le nombre de polygones. D'où
\[
  \Delta = \sum_{\gamma \in \Gamma_I} (-1)^{\nu_2(\gamma) + \nu_*^+(\gamma)}\,
  2^{\nu_1(\gamma) + \nu_*(\gamma)}\, M_\gamma, \qquad
  M_\gamma = \prod_{i \in I_1(\gamma)} Y_i \prod_{a \in A_2(\gamma)} Y_a^2
  \prod_{a \in A_*(\gamma)} Y_a ,
\]
où $\nu_1$ est le nombre de sommets isolés, $A_2$ l'ensemble des arêtes doubles et
$A_*$ celui des arêtes des polygones.

\textbf{Modulo $2$} (page 68). Seuls survivent les graphes sans sommet isolé ni
polygone, c'est-à-dire les involutions sans point fixe de $I$, et
\[
  \Delta \equiv \sum_{\gamma \in \Gamma_0} \Bigl(\prod_{a \in A(\gamma)} Y_a\Bigr)^2
  \pmod 2 .
\]
En particulier $\Delta \equiv 0 \pmod 2$ si $\mathrm{card}\, I$ est impair : c'est ce qui
rend le discriminant divisé $\delta' = \delta/2$ défini sur toute base en rang impair.

\textbf{Modulo $4$, en rang pair} $m = 2n$. Posons
\[
  B = \sum_{\gamma \in \Gamma_0} \prod_{a \in A(\gamma)} Y_a ,
\]
la somme sur les couplages parfaits de $I$ — l'hafnien de la matrice des $Y_{ij}$
— ; modulo $2$, c'est aussi le pfaffien de la matrice alternée à laquelle se réduit
la matrice de la forme, de sorte que $\Delta \equiv B^2 \pmod 2$. La page cherche
$\varepsilon = \pm 1$ tel que $B^2 - \varepsilon\Delta \equiv 0 \pmod 4$, et
« prévoit » $\varepsilon = (-1)^n$. Modulo $4$ ne restent dans $\Delta$ que les
graphes de $\Gamma_0$, de coefficient $(-1)^n$, et ceux de $\Gamma_1$ : sans sommet
isolé, avec un seul polygone, nécessairement de longueur paire, de coefficient
$\pm 2$. Dans
\[
  B^2 = \sum_{\gamma \in \Gamma_0} M_\gamma + 2 \sum_{\{\gamma, \gamma'\}} \prod_{a \in
  A(\gamma)} Y_a \prod_{a' \in A(\gamma')} Y_{a'},
\]
la seconde somme portant sur les paires non ordonnées de couplages distincts, le
produit de deux couplages est le monôme d'un graphe dont les composantes sont des
arêtes doubles et des cycles alternés de longueur paire $\geq 4$, et chaque tel
graphe provient de $2^{\nu_* - 1}$ paires\footnote{La page somme sur les couples
$\gamma \neq \gamma'$ ordonnés, avec le facteur $2$ (un « $4$ » surchargé), ce qui
compte chaque paire deux fois ; sa phrase « un même monôme intervient $2^{\nu_*}$
fois » est juste pour les couples ordonnés.}. Donc $B^2 \equiv \sum_{\Gamma_0}
M_\gamma + 2\sum_{\Gamma_1} M_\gamma \pmod 4$, et comme $2 \cdot (\pm 1) \equiv 2
\pmod 4$,
\[
  B^2 - (-1)^n \Delta \equiv 0 \pmod 4 , \qquad\text{i.e.}\qquad
  (-1)^n \Delta = B^2 - 4C, \quad C \in AQ_I .
\]
La page passe par une fausse alarme — une ligne à l'encre brune, « $\equiv 2\sum
M_\gamma \not\equiv 0$ (4) ! », corrigée par les deux lignes suivantes : chaque terme
vaut $2(1 \pm 1)M_\gamma \equiv 0$\footnote{Vérifié par machine pour $m = 2, 4,
6$ ; le choix $\varepsilon = -(-1)^n$ échoue dans chaque cas. Pour $m = 2$, c'est la
formule classique $b^2 - 4ac$.}. Le signe $(-1)^n$ est celui du discriminant
« changé de signe » ; c'est la forme universelle de la congruence de
Stickelberger\footnote{Le nom n'est pas sur la page.}.

\pagerange{87}{87}
\subsection*{La vérification en rang $4$ (page 87)}

Une page barrée d'un trait vérifie le développement pour $\mathrm{card}\,I = 4$ :
quatre sommets isolés, coefficient $2^4$, un graphe ; deux sommets et une arête
double, $-2^2$, six graphes ; un sommet et un triangle, $2^2$, quatre graphes (huit
permutations) ; deux arêtes doubles disjointes, $1$, trois graphes ; un carré,
$-2$, trois graphes (six permutations) ; au total $24 = 4!$ permutations, et les
coefficients sont ceux du déterminant\footnote{Contrôlé par machine. Les dernières
lignes de la page récrivent la congruence modulo $4$ ; un exposant y est illisible.}.

\pagerange{73}{79}
\subsection*{Pour un anneau : discriminants de la forme $b^2 - 4c$ (pages 73 à 79)}

Nouveau départ, sans titre. Soit $A$ un anneau, $\Delta \in A$, $A_0 = A/2A$.

\textbf{Lemme} (page 73). Pour tout $b_0 \in A_0$ il existe un unique $\gamma(b_0) \in
A/4A$ tel que $b^2 \equiv \gamma(b_0) \pmod 4$ pour tout $b \in A$ relevant $b_0$ :
$(b + 2x)^2 = b^2 + 4(x^2 + bx)$. Pour qu'il existe $b, c \in A$ avec $\Delta = b^2 -
4c$, il faut et il suffit que a) $\Delta_0$ soit le carré d'un $b_0 \in A_0$ (unique si
$A_0$ est réduit) et b) $\gamma(b_0) = \Delta \bmod 4$ — si $A_0$ n'est pas réduit,
pour l'un des $b_0$.

\textbf{Proposition} (page 77). Supposons $2$ régulier dans $A$, $A_0$ réduit et
intégralement clos, et $\Delta_0$ régulier dans $A_0$. Soit $\nu = 2d + 1$ un entier
impair. S'il existe $b', c' \in A$ et $v \in A$, de réduction $v_0$ régulière dans
$A_0$, tels que
\[
  b'^2 - 4c' = \Delta^\nu v^2,
\]
alors il existe $b, c \in A$ tels que $b^2 - 4c = \Delta$ (la réciproque est claire).
En effet, posant $w = \Delta^d v$, on a $\Delta_0 w_0^2 = b_0'^2$, donc
$(b'_0/w_0)^2 = \Delta_0$ dans l'anneau total des fractions de $A_0$ ; $A_0$ étant
intégralement clos, $b'_0 = b_0 w_0$ avec $b_0 \in A_0$, $b_0^2 = \Delta_0$. Si $b$
relève $b_0$, $bw$ relève $b'_0$, donc $b^2w^2 \equiv b'^2 \equiv \Delta w^2 \pmod
4$, et comme $w$ est régulier modulo $4$ ($2$ régulier dans $A$, $w_0$ régulier dans
$A_0$), $b^2 \equiv \Delta \pmod 4$\footnote{Les hypothèses (a) à (d) sont encadrées
sur les pages 73 à 77 ; la page énonce la Proposition avec « $v \in A^{*}$ » et
« $v_0$ $A_0$-régulier », et une note marginale précise qu'il suffit que $v$ soit tel
que $v_0$ soit régulier.}.

On considère alors $A_{2\Delta}$ et l'algèbre étale $B_{2\Delta} =
A_{2\Delta}[T]/(T^2 - \Delta)$, et l'on suppose qu'elle se prolonge en une algèbre
finie étale $B_\Delta$ sur $A_\Delta$. Si $\mathrm{Pic}\,A_\Delta = 0$ (e), $B_\Delta
\simeq A_\Delta[T]/(T^2 + bT + c)$, avec $b^2 - 4c \in A_\Delta^{*}$ puisque
$B_\Delta$ est étale, et, sur $A_{2\Delta}$, $b^2 - 4c = u^2\Delta$ avec $u \in
A_{2\Delta}^{*}$ ; si $A_\Delta$ est intégralement clos dans $A_{2\Delta}$ (f), $u
\in A_\Delta^{*}$ et $b^2 - 4c = \Delta u^2$ dans $A_\Delta$\footnote{La page conclut
« $u \in A^{*}$ » ; c'est $A_\Delta^{*}$ qu'on obtient. Un passage barré suppose
$(A_{2\Delta})^{*} = \{2^\mu\Delta^\nu v\}$, $v \in A^{*}$.}. La page 81 explique
comment chasser les dénominateurs : en multipliant $b$ et $c$ par $\Delta^\alpha$ et
$\Delta^{2\alpha}$, on se ramène à $b, c \in A$ et à une relation $b^2 - 4c =
\Delta^\nu u^2$ ; un passage barré montre qu'il faut pour cela connaître les unités
de $A_\Delta$, par exemple $A_\Delta^{*} = A^{*}\Delta^{\mathbf{Z}}$, pour retomber
sur la Proposition.

\pagerange{81}{86}
\subsection*{La Proposition de la page 85}

Au passage (page 81) : si $b^2 - 4c = \Delta u^2$ avec $u$ inversible,
$A_\Delta[T]/(T^2 + bT + c)$ est étale sur $A_\Delta$ et devient $A_{2\Delta}[U]/(U^2 -
\Delta)$ sur $A_{2\Delta}$, avec $U = (2T + b)/u$\footnote{La page écrit $U =
\frac{1}{2u}(T + \frac{b}{2})$, dont le carré est $\Delta/16$.}. La page 83, très
raturée, cherche à prouver qu'une unité $u_0$ est inversible, ébauche une première
Proposition aux hypothèses en partie illisibles, et note en marge qu'« il suffit
étale aux points maximaux de $\mathrm{Spec}\,A_0$ » ; la page 85 énonce au propre :

\textbf{Proposition} (page 85). Soit $A$ un anneau noethérien, $\Delta \in A$. On
suppose $2$ régulier dans $A$, $A_0 = A/2A$ intègre et normal, et $\Delta \notin 2A$.
Les conditions suivantes sont équivalentes :
\begin{itemize}
  \item[a)] il existe $b, c \in A$ avec $b^2 - 4c = \Delta$ ;
  \item[b)] il existe une extension quadratique $B$ de $A$ telle que $B_\Delta$ soit
    étale sur $A_\Delta$ et que $B_{2\Delta} \simeq A_{2\Delta}[U]/(U^2 - \Delta)$ ;
  \item[c)] il existe une extension finie étale $B_\Delta$ de $A_\Delta$ telle que
    $(B_\Delta)_2 \simeq A_{2\Delta}[U]/(U^2 - \Delta)$.
\end{itemize}
De plus, l'extension $B_\Delta$ de c) est déterminée à isomorphisme unique près.

Les implications a) $\Rightarrow$ b) $\Rightarrow$ c) sont claires, avec $B =
A[T]/(T^2 + bT + c)$. Pour c) $\Rightarrow$ a), on suit l'idée de la marge de la page
83 : localiser au point maximal de $\mathrm{Spec}\,A_0$. L'idéal $\mathfrak{p} = 2A$ est
premier, et $A_{\mathfrak{p}}$ est un anneau de valuation discrète d'uniformisante $2$ ;
comme $\Delta \notin \mathfrak{p}$, $A_{\mathfrak{p}}$ est un localisé de $A_\Delta$, et
$B_\Delta \otimes A_{\mathfrak{p}}$ est étale et libre de rang $2$, $\simeq
A_{\mathfrak{p}}[T]/(T^2 + \beta T + \gamma)$, avec $\beta^2 - 4\gamma = \Delta u^2$, $u
\in A_{\mathfrak{p}}^{*}$. Donc $\Delta_0$ est un carré dans le corps des fractions de
$A_0$, d'une racine $b_0$ entière, donc dans $A_0$ ; pour $b$ relevant $b_0$, $b^2 -
\Delta \in 4A_{\mathfrak{p}} \cap A = 4A$, la dernière égalité parce que $A_0$ est intègre
et $2$ régulier\footnote{Démonstration de l'édition ; les pages 77 à 84 suivent une
autre route, par $\mathrm{Pic}$ et les unités de $A_\Delta$, sous des hypothèses
supplémentaires. L'unicité de $B_\Delta$ est affirmée sans démonstration : celle de
l'isomorphisme est claire ($B_\Delta$ plat, $2$ régulier, donc $B_\Delta \subset
B_{2\Delta}$), son existence n'est pas établie sur les pages.}.

La page 86 esquisse la version globale : un schéma $X$, un diviseur $D$ (le lieu de
$\Delta$), $U = X \smallsetminus D$, $X_0 = V(2\mathcal{O}_X)$ intègre et normal, un
revêtement quadratique de $U$ dont on veut prolonger le module $L$ et le discriminant
; « OK sur $V = X - X_0$. Mais pas sur les points de $X_0$. » La page s'arrête sur
une accolade inachevée\footnote{La page 84 ne porte qu'une phrase : sous ces
conditions, il est nécessaire que l'extension soit étale aux points de $A_\Delta$.}.

\pagerange{88}{125}
\section*{V. Quadriques affines (pages 88 à 125)}

La chemise rose de la page 88 porte « Quadriques affines », suivi d'une suite
biffée : « et invariants des polyèdres réguliers sphériques ».

\pagerange{89}{97}
\subsection*{1. La sphère comme restriction de Weil d'une conique (pages 89 à 97)}

« “Invariants” des groupes des polyèdres réguliers sphériques. » Soit $(E, q,
\omega)$ un Module spécial quadratique de rang $3$ :
\[
  (1) \qquad \delta'(q) = -\,\omega^{\otimes -2}, \qquad \omega \text{ une base de }
  \det E
\]
\footnote{La page écrit « $\delta(q) = -\omega^{\otimes 2}$ ($\omega \in \Gamma(S, (\det
E)^{*})$) ». La convention de la page 105 (« $\omega_0^{\otimes -2} =
-\delta'(q)$ — attention au signe ! ») et celle de la page 128, (2.3), avec $\varepsilon
= -1$, sont celles qu'on suit ; c'est aussi celle du dossier 82.}. Il existe une
algèbre de quaternions tordue $M$, unique à isomorphisme unique près, et un
isomorphisme $E \simeq \gamma M = \mathrm{Ker}(M \xrightarrow{\mathrm{Trd}} \mathcal{O}_S)$
transformant $q$ en $q_M(u) = -\det u$ et $\omega$ en la base $\omega_M$ définie par
\[
  (4) \qquad \varphi_M([u, v], w)\, \omega_M = 2\, u \wedge v \wedge w ,
\]
où $\varphi_M(u, v) = \mathrm{Tr}(uv)$ est la polaire de $q_M$ sur $\gamma M$. Comme
$uv + vu = \mathrm{Tr}(uv)$ pour $u, v$ de trace nulle, $\mathrm{Tr}([u, v]w) =
2\mathrm{Tr}(uvw)$, et (4) s'écrit aussi
\[
  (8) \qquad u \wedge v \wedge w = \mathrm{Tr}(uvw)\, \omega_M
\]
\footnote{La lecture du dossier 82, à propos de ses formules (52) et (64), note
que cette normalisation donne la base $\omega$ opposée à la sienne. Une note
marginale de la page 89 dit $\omega$ « défini à $\pm$ près par (4) ».}. La loi
$[\ ,\ ]$ sur $E$ est transportée : $2\, x \wedge y \wedge z = \varphi([x, y], z)\,
\omega$.

On s'intéresse à la \emph{quadrique affine}
\[
  (10) \qquad \Sigma \subset \check{\mathbb{V}}(E), \qquad q(x) = -1,
\]
qui correspond à $\det u = 1$ dans $\gamma M$ : ce sont les $u$ de trace nulle tels que
$u^2 + 1 = 0$, c'est-à-dire les plongements de $\mathcal{O}_S[\mathbf{i}] =
\mathcal{O}_S[T]/(T^2 + 1)$ dans $M$. Soit $X$ la conique $q = 0$, $\tilde{X} \subset
\check{\mathbb{V}}(E)$ son cône épointé ($X = \tilde{X}/\mathbb{G}_m$), $S[\mathbf{i}] =
\mathrm{Spec}\,\mathcal{O}_S[\mathbf{i}]$, et $X^{S[\mathbf{i}]} =
\underline{\mathrm{Hom}}_S(S[\mathbf{i}], X)$ la restriction de Weil de
$X_{S[\mathbf{i}]}$ à $S$.

Une section de $\tilde{X}^{S[\mathbf{i}]}$ est un $\xi = x + \mathbf{i}y$, $x, y \in E$,
partout non nul, avec $q(x + \mathbf{i}y) = q(x) - q(y) + \mathbf{i}\varphi(x, y) = 0$ :
\[
  (17) \qquad q(x) = q(y), \qquad \varphi(x, y) = 0 .
\]
Le groupe $\mathbb{G}_m^{S[\mathbf{i}]}$ des $\alpha = a + \mathbf{i}b$ avec $a^2 + b^2$
inversible y opère, et l'ouvert $\tilde{U}$ où $q(x) = q(y)$ est inversible est stable :
$q(ax - by) = (a^2 + b^2)\, q(x)$\footnote{La page écrit le terme croisé $-2ab\,
\varphi(x, y)$ ; avec la convention $q(x + y) = q(x) + q(y) + \varphi(x, y)$ c'est
$-ab\,\varphi(x, y)$. Il est nul ici.}. Il définit un ouvert $U =
\tilde{U}/\mathbb{G}_m^{S[\mathbf{i}]} \subset X^{S[\mathbf{i}]}$. Pour $x, y \in E$ on a
\[
  (21\,\text{bis}) \qquad q([x, y]) = -4\, q(x)\, q(y) + \varphi(x, y)^2
\]
\footnote{Vérifié pour $\gamma M$ des matrices $2 \times 2$ de trace nulle. La page
l'écrit d'abord $-4q(x)q(y)\sin^2(\widehat{x, y})$ avec $\sin^2 = 1 - \frac{1}{4}
\varphi^2/q(x)q(y)$, la marge avertissant : « attention au signe et au facteur
$4$ ! »}, de sorte que, sur $\tilde{U}$, $q([x, y]) = -4q(x)^2$, et, si $2$ est
inversible,
\[
  (22) \qquad \tilde{f}(x + \mathbf{i}y) = [x, y]/2q(x) \in \Sigma .
\]
Comme $[ax - by, bx + ay] = (a^2 + b^2)[x, y]$, $\tilde{f}$ passe au quotient en $f : U
\to \Sigma$ (23).

\textbf{Proposition} (pages 94 à 97). Si $2$ est inversible, $f : U \to \Sigma$ est un
isomorphisme.

\emph{Monomorphisme.} Si $\tilde{f}(\xi) = \tilde{f}(\xi') = t$, alors $x, y$ ne sont
colinéaires sur aucune fibre (sinon $q(x) = 0$), sont orthogonaux à $t$, et $t$ n'est
isotrope sur aucune fibre ($\varphi(t, t) = 2q(t) = -2$) ; donc $(x, y)$ et $(x', y')$
sont deux bases de $t^{\perp}$. Les conditions $q(x') = q(y')$ et $\varphi(x', y') =
0$ forcent $x' = ax - by$, $y' = \lambda(bx + ay)$ avec $\lambda^2 = 1$, et l'égalité
des $[x', y']/2q(x')$ force $\lambda = 1$ : $\xi' = (a + \mathbf{i}b)\xi$.
\emph{Épimorphisme.} Pour $t \in \Sigma$, $t^{\perp}$ est un plan supplémentaire de
$\mathcal{O}_S t$ où $q$ est non dégénérée ; localement pour la topologie étale il a une
base orthonormée $(x, y)$, et $\tilde{f}(x + \mathbf{i}\lambda y) = t$ pour le signe
convenable $\lambda = \pm 1$.

\textbf{Corollaire.} Si $2$ est inversible, $g = f^{-1}$ est une immersion ouverte
\[
  (24) \qquad \Sigma \overset{g}{\hookrightarrow} X^{S[\mathbf{i}]} .
\]
Géométriquement : la polaire d'un point $t$ de la sphère coupe la conique en deux
points, conjugués au-dessus de $S[\mathbf{i}]$, et ce couple est un point de
$X^{S[\mathbf{i}]}$\footnote{Cette lecture géométrique est de l'édition ; c'est la
version, sur une base où $2$ est inversible, de l'identification de la sphère réelle
à la droite projective complexe, $\mathbb{P}^1(\mathbf{C}) = \mathrm{R}_{\mathbf{C}/\mathbf{R}}
(\text{conique})(\mathbf{R}) \smallsetminus \text{conique}(\mathbf{R})$.}.

\pagerange{97}{101}
\subsection*{Le complémentaire de $U$, et les essais de normalisation (pages 97 à 101)}

L'immersion canonique $X \hookrightarrow X^{S[\mathbf{i}]}$ (les points « réels »)
s'explicite par $x \mapsto x + \mathbf{i}\,0$.

\textbf{Proposition} (page 98). L'isomorphisme inverse de $f$ induit
\[
  \Sigma \xrightarrow{\ \sim\ } X^{S[\mathbf{i}]} \smallsetminus X = U
\]
\footnote{La page encadre « $\Sigma \xrightarrow{\sim} \tilde{X}^{S[\mathbf{i}]} -
\tilde{X} = \mathcal{U}$ », avec des tildes ; l'énoncé est au niveau des quotients par
$\mathbb{G}_m^{S[\mathbf{i}]}$. Une première version du lemme, barrée, demandait $\xi
\in \Gamma\tilde{\mathcal{U}}$ là où la conclusion vise le complémentaire.}, ce qui
équivaut au \textbf{Lemme} : sur un corps algébriquement clos de caractéristique
$\neq 2$, un $x + \mathbf{i}y$ avec $q(x) = q(y)$, $\varphi(x, y) = 0$, est hors de
$\tilde{U}$ (c'est-à-dire $q(x) = q(y) = 0$) si et seulement s'il est de la forme $(a +
\mathbf{i}b)x_0$ avec $q(x_0) = 0$ et $a^2 + b^2 \neq 0$. En effet $x$ et $y$
engendrent alors un sous-espace totalement isotrope d'une forme régulière de rang
$3$, donc $y = \lambda x$ avec $\lambda^2 \neq -1$, et $x + \mathbf{i}y = (1 +
\mathbf{i}\lambda)x$ ; la réciproque est immédiate. « O.K.… »

La suite de la page 99 essaie la caractéristique $2$ — $\mathbf{i} = 1 + \varepsilon$,
$X^{S[\mathbf{i}]}$ devient le fibré tangent à $X$ — et est barrée. La page 97 (barrée)
puis la page 100 (dans un cadre abandonné) cherchent à changer d'échelle pour que la
formule (21) prenne la forme $q'([x, y]') = q'(x)q'(y)\sin^2$ : avec $q_\lambda =
\lambda q$, $\omega_\mu = \mu\omega$, $[\ ,\ ]_{\lambda\mu} = (\lambda\mu)^{-1}[\ ,\ ]$,
il faut $\lambda^{-3}\mu^{-2} = -\tfrac{1}{4}$, ce que donnent $\lambda = -1$, $\mu =
\tfrac{1}{2}$, d'où $q' = -q$, $\omega' = \tfrac{1}{2}\omega$, $[x, y]' =
-2[x, y]$\footnote{La page écrit $\lambda^{-3}\mu^{-2} = -2^{-(3\alpha + 2\beta)} =
-2^{2}$ ; la condition $-\frac{1}{4}$ demande $-2^{-2}$, et c'est elle que résout
$3\alpha + 2\beta = -2$. La formule (26) s'interrompt après « $q''($ ».}. La page 101
s'ouvre sur un fragment barré qui fait passer $\tilde{f}$ au quotient sous le nom
$g : \tilde{X} \to \Sigma_{-4}$, puis note qu'en caractéristique $2$, $[x, y] =
\varphi(x, y)H$ où $H$ engendre le centre de $\tilde{\mathfrak{g}}$, de sorte que
l'involution s'y spécialise en l'involution constante : « rien à faire… ! »

\pagerange{102}{110}
\subsection*{1'. La version tordue : $X^{\Omega} \simeq \hat{\Sigma}$ (pages 102 à 110)}

Soit $E$ un fibré affine de rang $3$ et $\Sigma \subset E$ une quadrique affine
projectivement lisse. C'est la même chose qu'un fibré projectif $P$ de rang $3$ (la
clôture $\hat{E}$), une quadrique lisse $Q = \hat{\Sigma} \subset P$ et un plan $H$
(à l'infini) — ou encore une forme $Q$ de $\mathbb{P}^1 \times \mathbb{P}^1$ et un diviseur
de type $(1, 1)$ sur $Q$. On suppose que $\Sigma$ n'est un paraboloïde sur aucune
fibre : $H$ n'est pas tangent à $Q$, ou $X = H \cap Q$ est une conique lisse. En
termes d'extension $0 \to V \to \mathcal{E} \to \mathcal{O}_S \to 0$ ($V$ le fibré des
translations, de rang $3$) et d'une forme $Q$ sur $\mathcal{E}$ définie à une unité près,
la condition est que $q = Q|V$ soit lisse\footnote{La page écrit la suite exacte avec
un dernier terme « $1$ », pour $0$.}. En caractéristiques résiduelles $\neq 2$,
$\mathcal{E} = V \oplus V^{\perp}$, et $E$ a une origine canonique, le \emph{centre} $o$
de la quadrique, pôle de $H$ ; en normalisant $Q$ par $Q(o) = 1$, la donnée de
$\Sigma$ équivaut à celle du centre et d'une forme lisse $q$ sur $V$, et $\Sigma =
\{x \mid q(x) = -1\}$.

Soit alors $\Omega$ le sous-schéma de $\check{\mathbb{V}}(\det V^{\vee})$ des $\omega$
tels que $\omega^{\otimes -2} = \delta'(q)$ : un $\mu_2$-torseur, c'est-à-dire, en
caractéristiques résiduelles $\neq 2$, un revêtement étale de degré $2$. Si $q$ est
spécial orthogonal, muni de $\omega_0$ avec $\omega_0^{\otimes -2} = -\delta'(q)$, les
$\omega$ sont les $\lambda\omega_0$ avec $\lambda^2 = -1$, et $\Omega =
S[\mathbf{i}]$. La symétrie $\sigma$ de centre $o$ a pour points fixes dans $\hat{E}$
la section $o$ et le plan $H$ ; dans $\hat{\Sigma}$, la conique $X$.

\textbf{Le programme} (page 106). L'automorphisme canonique de $\Omega$ définit un
automorphisme $\sigma$ de $X^{\Omega} = \underline{\mathrm{Hom}}_S(\Omega, X)$, et
$(X^{\Omega})^{\sigma} \simeq X$. On veut un isomorphisme canonique $X^{\Omega}
\xrightarrow{\sim} \hat{\Sigma}$, compatible à $\sigma$ et prolongeant l'inclusion
de $X$, d'où $\Sigma \simeq X^{\Omega} \smallsetminus X$. Par descente le long de
$\Omega$, puisque $(X^{\Omega})_{\Omega} \simeq X_{\Omega} \times X_{\Omega}$, il
revient au même de construire, pour toute section $\omega$ de $\Omega$ (sur une base
$S'$ quelconque), un isomorphisme $f_\omega : X \times X \to \hat{\Sigma}$
satisfaisant
\begin{itemize}
  \item[a)] la compatibilité aux changements de base ;
  \item[b)] $f_{-\omega}(x, y) = f_\omega(y, x)$ ;
  \item[c)] $f_\omega(y, x) = \sigma f_\omega(x, y)$ ;
  \item[d)] $f_\omega(x, x) = i(x)$, $i : X \hookrightarrow \hat{\Sigma}$.
\end{itemize}
Pour $\omega$ donné, $(E, -q, \omega)$ est spécial pour la convention de la page 89,
d'où une algèbre de quaternions $M$ avec $(E, -q, \omega) \simeq (\gamma M, q_M,
\omega_M)$, donc $q = \det$ sur $\gamma M$. Alors $X$ est la conique des éléments
nilpotents ($\mathrm{Tr} = \det = 0$), et $\Sigma$ ($\mathrm{Tr}\,u = 0$, $\det u =
-1$) le schéma des \emph{réflexions} $u$ ($u^2 = 1$) associées à $M$\footnote{La page
écrit « transformant $-q$ en $q_M = -\det|\gamma M$ », avec en marge « attention au
signe » ; c'est bien $-q = -\det$, soit $q = \det$.}. Comme $M = \gamma M \oplus
\mathcal{O}_S$ ($2$ inversible), $\hat{\Sigma}_M \subset \check{\mathbb{P}}(M)$ a pour
équation $\det(u + t\cdot 1) = \det u + t^2 = 0$ : c'est la quadrique $\det v = 0$
des endomorphismes de rang $\leq 1$, et l'on « sait » l'isomorphisme canonique
$F_M : X_M \times X_M \simeq \hat{\Sigma}_M$ — à un couple de droites $(\ell, \ell')$
d'un module de rang $2$, l'endomorphisme de rang $1$ de noyau $\ell$ et d'image $\ell'$
— qui induit sur la diagonale l'inclusion de $X$, et qui transforme la symétrie
des facteurs en l'involution $v \mapsto v' = \mathrm{Tr}(v)\cdot 1 - v$ (l'adjointe,
qui conserve le déterminant). Pour $v = u + t\cdot 1$, $-v' = u - t\cdot 1$ : sur
$\hat{\Sigma}_M$, c'est l'antipodie $\sigma$. D'où a), c), d) ; b) résulte de la
fonctorialité en $(E, q, \omega)$ appliquée à $-\mathrm{id}_E : (E, q, \omega) \to
(E, q, -\omega)$ (en rang $3$, $\det(-\mathrm{id}) = -1$), qui induit $\sigma$ sur
$\hat{\Sigma}$.

\pagerange{110}{116}
\subsection*{La jonction avec le § 1, et le signe (pages 110 à 116)}

Reste à comparer, quand $\omega_0$ est donné avec $\delta'(q) = -\omega_0^{\otimes -2}$
(le cadre du § 1, $\Omega = S[\mathbf{i}]$), l'isomorphisme $f_{q, \omega} :
X^{\Omega} \smallsetminus X \simeq \Sigma$ et le $g_q$ du § 1. On a $f_{q, \omega} =
\theta_{q, \omega}\, g_q$ pour un automorphisme $\theta_{q, \omega}$ de $\Sigma$,
fonctoriel et compatible aux changements de base. Sur $S = \mathrm{Spec}\,
\mathbf{Z}[\tfrac{1}{2}]$ et pour la forme standard, $\Sigma \simeq G/T$ avec $G =
\mathrm{SO}(3)$ et $T$ le tore maximal stabilisateur d'un point ; les automorphismes
de $G/T$ commutant à $G$ forment $N(T)/T \simeq \mathbf{Z}/2$, engendré par
l'antipodie. « Mais lequel des deux ? »\footnote{Une tentative, encadrée et barrée
page 111, passe par $\theta : x \mapsto ix$ de $\gamma M$ dans lui-même, qui transforme
$q_M$ en $-q_M$ et $\omega_M$ en $i^3\omega_M$ ; la page note « TSVP ».}

Le calcul (pages 113 et 114) : avec $(E, q, \omega_0) = (\gamma M, q_M, \omega_M)$ et
$\theta = i\,\mathrm{id}$, $(u + \mathbf{i}v, u - \mathbf{i}v)$ devient $(\alpha,
\beta) = (iu - v, iu + v)$, d'où $2i[u, v] = [\alpha, \beta]$ et $-4q(u) =
\mathrm{tr}(\alpha\beta)$, et l'on trouve que l'isomorphisme cherché envoie un couple
$(\alpha, \beta)$ d'éléments nilpotents, « assez disjoints » ($\mathrm{tr}(\alpha\beta)$
inversible), sur
\[
  -[\alpha, \beta]/\mathrm{tr}(\alpha\beta) \in \Sigma
\]
— « et non $+[\alpha, \beta]/\mathrm{tr}(\alpha\beta)$ » —, c'est-à-dire, dans
$\hat{\Sigma}$, sur le point de coordonnées $\mathrm{tr}(\alpha\beta)\cdot 1 - [\alpha,
\beta]$. Pour $\alpha, \beta$ nilpotents de $\gamma M$, $\alpha\beta + \beta\alpha =
\mathrm{tr}(\alpha\beta)$, de sorte que
\[
  F_M(\alpha, \beta) = \mathrm{tr}(\alpha\beta)\cdot 1 - [\alpha, \beta] = 2\beta\alpha,
\]
l'endomorphisme de rang $1$ de noyau la droite de $\alpha$ et d'image celle de
$\beta$\footnote{La page présente $F_M$ comme une application « bilinéaire au niveau
des cônes affines » $C_M \times C_M \to D_M$, compatible à $(\lambda, \mu) \mapsto
\lambda\mu$, induisant l'isomorphisme $X_M \times X_M \simeq \hat{\Sigma}_M$. Mais
$F_M(\alpha, \alpha) = 2\alpha^2 = 0$ : la formule n'est définie que hors de la
diagonale, où elle donne bien l'isomorphisme ; sur la diagonale, l'isomorphisme
envoie $\alpha$ sur $\alpha$ lui-même, et ne provient pas de cette application
bilinéaire. L'identité $2\beta\alpha$ est de l'édition.}. Qu'il soit dans
$\hat{\Sigma}_M$ résulte de $q_M([\alpha, \beta]) = \varphi(\alpha, \beta)^2$, cas
particulier de (21 bis).

La page 116 généralise cette identité à deux éléments quelconques $u = \alpha + s$,
$v = \beta + t$ de $M$, avec $\mathrm{tr}\,u = 2s$, $\mathrm{tr}\,v = 2t$ : on a $[u, v] =
[\alpha, \beta]$, $\mathrm{tr}(\alpha\beta) = \mathrm{tr}(uv) - \tfrac{1}{2}
\mathrm{tr}\,u\,\mathrm{tr}\,v$, $\det\alpha = \det u - \tfrac{1}{4}(\mathrm{tr}\,u)^2$, et
l'identité « encadrée » devient
\[
  \det[u, v] + (\det u)(\mathrm{tr}\,v)^2 + (\det v)(\mathrm{tr}\,u)^2 - 4(\det u)(\det
  v) + (\mathrm{tr}\,uv)^2 - (\mathrm{tr}\,uv)(\mathrm{tr}\,u)(\mathrm{tr}\,v) = 0,
\]
valable pour toutes matrices $2 \times 2$ sur tout anneau\footnote{La page écrit
$\det(u - s) = \det(s - u) = s^2 + \det u$ ; pour une matrice $2 \times 2$, $\det(u -
s\cdot 1) = \det u - s\,\mathrm{tr}\,u + s^2 = \det u - s^2$ ici. L'erreur passe dans
l'encadré, qui porte $-(\det u)(\mathrm{tr}\,v)^2 - (\det v)(\mathrm{tr}\,u)^2$ : sous
cette forme l'identité est fausse, et elle est vraie avec les signes $+$ (vérifié par
machine, sur $\mathbf{Z}$).}.

\pagerange{117}{121}
\subsection*{1''. Quadriques affines birégulières (pages 117 à 121)}

« Quadriques affines birégulières (i.e.\ lisses et lisses à l'infini,
projectivement). » Soient $P$ un fibré projectif, $Q \subset P$ une quadrique lisse,
$H$ un hyperplan non tangent à $Q$, $E = P \smallsetminus H$, défini par une extension
\[
  (1) \qquad 0 \to V \to T \xrightarrow{\ \pi\ } \mathcal{O}_S \to 0, \qquad E =
  \pi^{-1}(1),
\]
et $q$ une forme sur $T$, définie à une unité près, qui définit $Q$. $Q$ lisse veut
dire $q$ régulière ; $H$ non tangent veut dire $q|V$ régulière.

\textbf{$T$ de rang pair $2\nu$} (pages 117 à 120). La forme polaire est alors non
dégénérée en toute caractéristique, et $\mathbb{D} = V^{\perp}$ est un sous-Module
inversible de $T$ ; $H$ est tangent à $Q$ au point $p$ exactement quand $H = p^{\perp}$
avec $q(p) = 0$, donc $H$ est non tangent si et seulement si $q|\mathbb{D}$ ne s'annule
sur aucune fibre — ce que la marge demandait de « vérifier ».

\textbf{Proposition.} Il existe une unique section $a$ de $\mathbb{D}$ et une unique
représentante $q$ de la forme telles que
\[
  (8) \qquad q(a) = 1, \qquad (9) \qquad \pi(x) = \varphi(a, x) \quad (x \in T) ;
\]
alors $a$ est une base de $\mathbb{D}$, et
\[
  (10) \qquad \pi(a) = 2 .
\]
(Si $(a_0, q_0)$ est un choix quelconque, $\varphi_0(a_0, \cdot) = \mu\pi$ avec $\mu$
inversible, et $(\lambda a_0, \kappa q_0)$ convient pour $\lambda = \mu/q_0(a_0)$,
$\kappa = 1/\lambda\mu$, et pour eux seuls ; (10) est (9) en $x = a$.)\footnote{La
description de $E$ dans (9) est biffée par hachures sur la page, ne laissant que
« $\varphi(a, x)$ », et la formule (11) $\pi(x) = \varphi(a, x)$ de la page 120 est
biffée ; c'est pourtant elle qui rend (10) vraie et l'unicité exacte. On l'adopte.}

\textbf{Scholie} (page 120). La catégorie des $(P, Q, H)$ équivaut à celle des $(T, q,
a)$ : $T$ de rang $2\nu$, $q$ régulière, $a$ une section de $T$ avec $q(a) = 1$ ; on
reprend $\pi = \varphi(a, \cdot)$ et $V = a^{\perp}$. (En caractéristique $2$, $a \in
a^{\perp}$ est le radical de $q|V$, sur lequel $q$ ne s'annule pas : $q|V$ est encore
régulière.)

\emph{N.B.} a) En caractéristiques résiduelles $\neq 2$, $T = V \oplus \mathbb{D}$
orthogonalement, $q(x + ta) = q_0(x) + t^2$, et $E = \{t = \tfrac{1}{2}\}$ : $Q \cap E$
a pour équation $q_0(x) = -\tfrac{1}{4}$, soit $q_0(x) = -1$ après avoir remplacé $q_0$
par $4q_0$. La catégorie des $(E, Q \cap E)$ équivaut alors à celle des $(V, q_0)$, $V$
de rang $2\nu - 1$\footnote{La page écrit $q(x, t) = q_0(x) + t^2$ et « $q_0(x) + 1 =
0$ », ce qui suppose $\pi(a) = 1$, la normalisation de la page 121 en rang impair ;
avec (10), c'est $t = \frac{1}{2}$.}. En caractéristique $2$, au contraire, $\mathbb{D}
\subset V$, et il n'y a pas de décomposition canonique.

\textbf{$T$ de rang impair $2\nu + 1$} (pages 120 et 121). Alors $V$ est de rang pair
et $q|V$ régulière a une polaire non dégénérée en toute caractéristique, de sorte que
$T = V \oplus \mathbb{D}$ orthogonalement, $\mathbb{D} \to T/V \simeq \mathcal{O}_S$ étant un
isomorphisme (14). D'où une section canonique $a$ de $\mathbb{D}$ avec $\pi(a) = 1$
(15) ; on normalise $q(a) = 1$ (16), et
\[
  (17)\text{--}(18) \qquad q(x, t) = q_0(x) + t^2, \qquad Q \cap E : \ q_0(x) = -1,
\]
« sans restriction à la car.\ $\neq 2$ » : les quadriques affines birégulières, dans un
fibré affine de rang pair, équivalent aux formes régulières $q_0$ sur le fibré des
translations.

\pagerange{122}{125}
\subsection*{2--3. La symétrie canonique et le groupe d'automorphismes (pages 122 à 125)}

Une réflexion de $P$ d'hyperplan fixe $H$ est, à un scalaire près, $\sigma = \mathrm{id} +
\pi \otimes c$, $c \in T$. Elle préserve $q$ à une unité $\lambda$ près si $(\lambda -
1)q = \pi\,(q(c)\pi + \varphi(c, \cdot))$ ; si $\mathrm{rg}\,T \geq 3$, cela force
$\lambda = 1$ et
\[
  (19) \qquad q(c)\,\pi + \varphi(c, \cdot) = 0 .
\]
En rang pair, où $\pi = \varphi(a, \cdot)$, (19) donne $c = \lambda a$ avec $\lambda
= -\lambda^2$, d'où (hors de l'identité) $c = -a$ :
\[
  (20) \qquad \sigma(x) = x - \pi(x)\, a, \qquad (21) \qquad \sigma^2 = \mathrm{id},
\]
puisque $\pi(a) = 2$\footnote{La page justifie (21) par « $\pi(a) = 1$ » ; avec $\pi(a)
= 1$, $\sigma^2 = \sigma$. Elle étiquette (20) « rg $E$ pair » après avoir dit
s'intéresser surtout au cas « $E$ de rg impair » : il s'agit de $T$ de rang pair, $E$
de rang impair.}. \textbf{Proposition} (en marge) : il existe une unique réflexion de
$P$ d'hyperplan fixe $H$ laissant $Q$ invariante, $\sigma = \mathrm{id} - \pi \otimes a$ :
\emph{la symétrie canonique} de la quadrique affine. En caractéristique $\neq 2$ elle
induit sur $E$ la symétrie de centre $a/2$ ; en caractéristique $2$, $a \in V$ et c'est
une transvection.

Un feuillet découpé (page 124) donne l'exemple $T = M_2$, $q = \det$, $a = 1$ : la
polaire de $\det$ est $\varphi(x, x') = ad' + da' - bc' - cb'$, d'où $\pi(x) = \varphi(x,
1) = \mathrm{Tr}\,x$, $\pi(a) = 2$, et $V = \gamma M$, les matrices de trace
nulle\footnote{La page écrit $\varphi = bc' + cb' - (ad' + da')$, l'opposé, d'où $\pi =
-\mathrm{Tr}$ et la remarque « attention au signe ! ». Un fragment tête-bêche, coupé
par la découpe, n'est pas utilisé.}. En rang impair (page 125), $\sigma(x) = x -
2\pi(x)\,a$ (22) est la symétrie par rapport à l'origine.

\textbf{§ 3.} Soit $G = \underline{\mathrm{Aut}}(P, H, Q) =
\underline{\mathrm{Aut}}_{\mathrm{aff}}(E, Q \cap E)$. Si $T$ est de rang impair, $G
\simeq \underline{\mathrm{Aut}}(V, q_0)$, $V$ de rang pair, extension de $\mathbf{Z}/2$
par le sous-groupe $G^{+}$ (la page le définit comme opérant trivialement sur $\det V$
; en caractéristique $2$ il faut l'invariant de Dickson de la partie VI). Si $T$ est de
rang pair, « on va montrer » que $G \simeq G^0 \times (\mathbf{Z}/2)_S$, avec $G^0
\simeq \underline{\mathrm{Aut}}^{+}(V, q_0)$ et $\mathbf{Z}/2$ engendré par $\sigma$ ; la
page s'arrête là.

\section*{VI. Formes quadratiques, fibrés en drapeaux, groupes orthogonaux (pages 126 à 157)}

La feuille de garde de la page 126 porte « Formes quadratiques. Fibrés en drapeaux
quad., groupes orthogonaux » et « (31~p) ».

\pagerange{127}{132}
\subsection*{1--4. Discriminant, régularité, groupes orthogonaux (pages 127 à 132)}

\textbf{1.} Pour $q$ de rang $n$ sur $E$, le discriminant divisé $\delta'(q) \in
\Gamma(\det E)^{\otimes -2}$ vérifie $\delta'(\lambda q) = \lambda^n\delta'(q)$. Sont
équivalentes : a) $\delta'(q)$ inversible ; b) $V(q) \subset \check{\mathbb{P}}(E)$ est
un diviseur relatif lisse ; c) en chaque point géométrique, $q$ est dans l'orbite
ouverte de $\mathrm{GL}$ sur les formes, de dimension $n(n+1)/2$ ; d) le groupe
d'automorphismes de chaque fibre géométrique est de dimension minimale $n(n-1)/2$ ;
e) $q$ est localement isomorphe à la forme standard $q_n$ — pour la topologie
étale si $2$ est inversible ou $n$ pair, pour la topologie fidèlement plate en
général. Pour une forme à valeurs dans un Module inversible $L$, $\delta'(q)$ est une
section de $(\det E)^{-2} \otimes L^{\otimes n}$ et, $q$ étant régulière, un
isomorphisme $L^{\otimes n} \simeq (\det E)^{\otimes 2}$ (1.3).

\textbf{2. Rang impair $n = 2\nu + 1$.} $q_n = \sum_{i} x_iy_i + z^2$,
$\delta'(q_n) = (-1)^\nu$. On a $\mathbb{G}_m \cap O(q) = \mu_2 =
\underline{\mathrm{Centre}}\,O(q)$, et $\mu_2 \cap \mathrm{SO}(q) = 1$.
\textbf{Théorème} : $O(q) = \mu_2 \times \mathrm{SO}(q)$, et $\mathrm{SO}(q)$ est lisse à
fibres connexes, semi-simple (absolument presque simple) si $n \geq 3$\footnote{La
page ajoute « simple si $n \geq 7$ » ; en rang impair $\mathrm{SO}_{2\nu+1}$ est de type
$B_\nu$, simple pour tout $\nu \geq 1$. La restriction vaudrait en rang pair, où
$\mathrm{SO}_4$ n'est pas simple.}. Une \emph{structure spéciale orthogonale d'indice
$\varepsilon$} ($\varepsilon = \pm 1$) est une base $\omega$ de $\det E$ avec
\[
  (2.3) \qquad \omega^{\otimes -2} = \varepsilon\,\delta'(q) ;
\]
alors $\underline{\mathrm{Aut}}(E, q, \omega) = \mathrm{SO}(q)$ ; $q_n$ est un prototype
pour $\varepsilon = (-1)^\nu$, $-q_n$ pour $\varepsilon = (-1)^{\nu+1}$, et $(q, \omega)
\mapsto (-q, \omega)$ échange les deux indices.

\textbf{Proposition} (classification ; page 130, qui précède la page 129 dans
l'ordre de lecture). Soit $\varepsilon = \pm 1$. La catégorie des modules
quadratiques réguliers $(E, q)$ de rang $n = 2\nu + 1$ équivaut à celle des $(E_0,
q_0, \omega_0, \Delta, \phi)$, où $(E_0, q_0, \omega_0)$ est spécial orthogonal
d'indice $\varepsilon$ et $(\Delta, \phi : \Delta^{\otimes 2} \simeq \mathcal{O}_S)$
un $\mu_2$-torseur ; le foncteur est $E = E_0 \otimes \Delta$, $q = q_0 \otimes \phi$,
et le quasi-inverse $\Delta = \det E$, $\phi$ donné par $\delta'(q)$, $E_0 = E \otimes
\Delta^{-1}$, $q_0 = \varepsilon\,q \otimes \Delta^{-1}$ (via $\phi$), $\omega_0$ venant
de $\det E_0 \simeq \Delta^{1-n} = \Delta^{-2\nu} \simeq \mathcal{O}_S$\footnote{On a
vérifié, dans une base locale, que $\omega_0^{\otimes -2} = \varepsilon\,\delta'(q_0)$
avec ces choix. La page note le second facteur « $q_0 \otimes L$ ».}. NB :
$\underline{\mathrm{Aut}}\,V(q) \simeq O(q)/\mu_2 \simeq \mathrm{SO}(q)$, lisse, connexe,
réductif.

\textbf{3. Similitudes.} Pour $q : E \to L$, $O'(q) \subset \mathrm{GL}(E) \times
\mathbb{G}_m$ est le groupe des $(u, \lambda)$ tels que $q(ux) = \lambda q(x)$ ; il
contient le centre $\mathbb{G}_m$, sur lequel le multiplicateur est $t \mapsto t^2$.
\textbf{Proposition} (page 129) : en rang impair, $O'(q) \simeq \mathrm{SO}(q) \times
\mathbb{G}_m$, par $(g, t) \mapsto (tg, t^2)$, d'inverse $t = \det u/\lambda^\nu$.
\textbf{Corollaire} : pour $\varepsilon$ fixé, les modules quasi-orthogonaux $(E, L, q)$
de rang impair équivalent aux $(E_0, q_0, \omega_0, \Delta)$, $\Delta$ inversible, par
$L = \Delta^{\otimes 2}$, $E = E_0 \otimes \Delta$, et réciproquement $\Delta =
L^{\otimes -\nu} \otimes \det E$\footnote{Le calcul de $\det E_0$ en bas de la page 129
porte $L$ là où il faut $L^n$, et plusieurs essais biffés.}.

\textbf{4. Rang pair $n = 2\nu$.} $q_{2\nu} = \sum x_iy_i$, $\delta'(q_{2\nu}) =
(-1)^\nu$, et $\mu_2 \subset O(q) \cap \mathrm{SL}(E)$ ; « on ne l'appelle pas
$\mathrm{SO}(q)$ à cause des car.\ $2$ ». \textbf{Théorème} : $O(q)$ est lisse,
$O(q)/\mathrm{SO}(q) \simeq (\mathbf{Z}/2)_S$ avec $\mathrm{SO}(q)$ sa composante neutre, et
$\mu_2 = \underline{\mathrm{Centre}}\,O(q) \subset \mathrm{SO}(q)$. \textbf{Corollaire} :
pour $X = V(q)$ et $n \geq 3$, $\underline{\mathrm{Aut}}(X) \simeq O(q)/\mu_2$ est lisse,
de composante neutre $\mathrm{SO}(q)/\mu_2$ et de groupe des composantes
$\mathbf{Z}/2$\footnote{La restriction $n \geq 3$ est de l'édition : en rang $2$, $X$ est
un couple de points.}. L'essentiel est l'épimorphisme $O(q) \to \mathbf{Z}/2$, défini par
un revêtement étale de degré $2$ attaché à $q$ : \textbf{Proposition} (page 132) : le
sous-schéma de $\underline{\mathrm{Grass}}_\nu(E)$ des sous-espaces totalement isotropes
maximaux est fermé, lisse, et ses fibres géométriques ont deux composantes connexes.

\pagerange{133}{141}
\subsection*{Fibrés en drapeaux isotropes (pages 133 à 141)}

Soit $(E, q)$ régulière de rang $N = 2n + 1$ ou $N = 2n$, $n \geq 1$, et $H = \{i_1 <
\cdots < i_\nu\} \subset [1, n]$. Soit $\underline{\Delta}_H(E, q)$ le foncteur des
drapeaux $E_1 \subset \cdots \subset E_\nu$ de type $H$ (les $E_k$ facteurs directs de
rangs $i_k$) avec $E_\nu$ totalement isotrope ($q|E_\nu = 0$).

\textbf{Théorème.} a) $\underline{\Delta}_H(E, q) \hookrightarrow
\underline{\mathrm{Drap}}_H(E)$ est une immersion fermée ; $\underline{\Delta}_H$ est donc
représentable et projectif sur $S$. b) Il est lisse sur $S$, à fibres géométriques
connexes, sauf si $N = 2n$ et $i_\nu = n$, auquel cas elles ont exactement deux
composantes connexes. (En marge : même énoncé pour $\underline{\Delta}_H \to
\underline{\Delta}_K$, $K \subset H$.)

\emph{Lissité.} Par le critère infinitésimal, il s'agit de relever, au-dessus d'un
$S$ affine et d'un idéal $J$ de carré nul, un sous-module totalement isotrope $F_0$ de
$E_0$. On relève $F_0$ en un facteur direct $F'$ quelconque ; $q|F'$ est alors à
valeurs dans $J$, et l'on corrige $F'$ par un $u \in \mathrm{Hom}(F_0, E_0/F_0) \otimes
J$ grâce à la surjectivité de
\[
  \underline{\mathrm{Hom}}(F_0, E_0/F_0) \to \underline{\mathrm{Hom}}(F_0,
  E_0/F_0^{\perp}) \simeq \underline{\mathrm{Hom}}(F_0, F_0^{\vee}) \to
  \underline{\mathrm{Quad}}(F_0), \qquad u_0 \mapsto (x \mapsto \varphi_0(x, u_0x))
\]
— toute forme quadratique sur un module projectif est de la forme $B(x, x)$ —, et du
\textbf{Lemme 1} : si $F$ est totalement isotrope et facteur direct, $E \to E^{\vee}
\to F^{\vee}$ est un épimorphisme, $F^{\perp} \supset F$ est facteur direct, de formation
compatible aux changements de base, et $F \times E/F^{\perp} \to \mathcal{O}_S$ est une
dualité. Sur un corps algébriquement clos, le seul cas non évident est celui d'une
forme dégénérée pour $\varphi$, en caractéristique $2$ et rang impair : un $a \in F$
orthogonal à tout $E$ et vérifiant $q(a) = 0$ est nul, $q$ étant régulière\footnote{Une
première tentative de relèvement, par récurrence sur le rang de $F$ et vecteur par
vecteur, est barrée (pages 134 et 135).}.

\emph{Connexité.} \textbf{Lemme 3} : pour $F$ totalement isotrope de rang $d$, $E' =
F^{\perp}/F$ est régulière de rang $N - 2d$, et les $G \supset F$ totalement isotropes
de rang $d'$ correspondent aux sous-espaces totalement isotropes de rang $d' - d$ de
$E'$. \textbf{Lemme 2} : les fibres de $\underline{\Delta}_{\{d, d'\}} \to
\underline{\Delta}_{\{d\}}$ sont lisses, projectives, à fibres géométriques connexes, sauf
si $N$ est pair et $d' = n$, où elles ont deux composantes. \textbf{Lemme 4} :
$\underline{\Delta}_{\{1\}}(E)$, la quadrique elle-même, a des fibres géométriques
connexes si $\mathrm{rg}\,E \geq 3$, et formées de deux points si $\mathrm{rg}\,E = 2$.
Par dévissage, cela règle tout sauf $N = 2n$, $H = \{n\}$.

\textbf{Le cas $N = 2n$, $H = \{n\}$} (pages 138 à 140). Pour $F$ totalement isotrope
de rang $n$, $F = F^{\perp}$, $E/F \simeq F^{\vee}$, d'où un isomorphisme canonique
$\det F \otimes \det F^{\vee} \simeq \mathcal{O}_S \to \det E$, c'est-à-dire une section
$\omega_F$ de $\det E$ ; si $(e_1, \ldots, e_n)$ est une base de $F$ et $(e_i^{*})$ la
base duale d'un supplémentaire $G \simeq F^{\vee}$, $\omega_F = e_1 \wedge \cdots
\wedge e_n \wedge e_1^{*} \wedge \cdots \wedge e_n^{*}$\footnote{La page écrit « base de
$E$ » ; il s'agit d'une base de $F$.}. \textbf{Lemme 5} : $F$ a un supplémentaire
totalement isotrope $G$, et alors $q(x \oplus y) = \langle x, y\rangle$. Le calcul du
déterminant donne $\delta(q) = (-1)^n\,\omega_F^{\otimes -2}$, soit
\[
  \omega_F^{\otimes 2} = (-1)^n\,\delta(q)^{-1}
\]
\footnote{La page écrit d'abord $\delta(q) = (-1)^n(\omega_F)^2$, puis encadre la forme
juste.}. D'où un morphisme $\underline{\Delta}_{\{n\}}(E, q) \to
\underline{\mathrm{Arf}}'(E, q)$, $F \mapsto \omega_F$, vers le schéma des racines carrées
de $(-1)^n\delta(q)^{-1}$ dans $\check{\mathbb{V}}(\det E)$ — un revêtement étale de
degré $2$ si $2$ est inversible. \textbf{Théorème 2} ($2$ inversible) : ce morphisme
est lisse, projectif, à fibres géométriques connexes. Il est surjectif parce que $O(q)$
opère transitivement sur $\underline{\mathrm{Arf}}'$ — une réflexion échange $\omega$
et $-\omega$ — ; la connexité des fibres se ramène aux drapeaux complets et au
\textbf{Lemme 6} (page 141) : $\underline{\Delta}^{\omega}_{\{n-1, n\}} \to
\underline{\Delta}_{\{n-1\}}$ est un isomorphisme, c'est-à-dire que tout $F'$ totalement
isotrope de rang $n - 1$ est contenu dans exactement un $F$ de rang $n$ avec $\omega_F =
\omega$. En effet les $F \supset F'$ correspondent aux deux droites isotropes du plan
hyperbolique $F'^{\perp}/F'$ (Lemme 3), et les deux $\omega_F$ sont opposés\footnote{La
page écrit « $F \subset F'$ » ; vu les rangs, c'est $F' \subset F$. Elle note en marge
que l'argument est « indépendant de la question de car. »}. Une fois traité le cas
d'un corps, on passe au cas général par spécialisation et factorisation de Stein ; et
c'est ainsi qu'on définit $\underline{\mathrm{Arf}}(E, q)$ en général, y compris en
caractéristique $2$.

\pagerange{142}{145}
\subsection*{Récapitulation, et le revêtement d'Arf (pages 142 à 145)}

\textbf{Théorème.} Soit $(E, q)$ régulière de rang $N = 2n$ ou $2n + 1$, $H' \subset H
\subset [1, n]$, $d = \max H$, $d' = \max H'$.
\begin{itemize}
  \item[a)] $p_{H', H} : \underline{\Delta}_H(E, q) \to \underline{\Delta}_{H'}(E, q)$ est
    projectif, lisse et surjectif ; en particulier ($H' = \varnothing$) les
    $\underline{\Delta}_H \to S$ le sont.
  \item[b)] Ses fibres géométriques sont connexes, sauf si $N$ est pair et $d' < d =
    n$, où elles ont deux composantes connexes.
  \item[c)] Pour $N = 2n$, le schéma $\underline{\mathrm{Arf}}(E, q)$ des composantes
    connexes de $\underline{\Delta}_{\{n\}}(E, q)$ sur $S$ est un revêtement étale de
    degré $2$ ; si $n \in H$, $\underline{\Delta}_H \to \underline{\Delta}_{\{n\}} \to
    \underline{\mathrm{Arf}}$ est projectif, lisse, à fibres géométriquement connexes :
    c'est la factorisation de Stein de $\underline{\Delta}_H \to S$. Plus généralement,
    si $d' < d = n$, $\underline{\Delta}_H \to \underline{\Delta}_{H'} \times_S
    \underline{\mathrm{Arf}}$ est à fibres géométriquement connexes.
  \item[d)] Pour $N = 2n$, $\underline{\Delta}_{\{n-1, n\}} \simeq
    \underline{\Delta}_{\{n-1\}} \times_S \underline{\mathrm{Arf}}$ ; en particulier
    $\underline{\Delta}_{\{n-1, n\}} \to \underline{\Delta}_{\{n-1\}}$ est fini étale de
    degré $2$.
  \item[e)] Soit $\underline{\mathrm{Arf}}' \subset \check{\mathbb{V}}(\det E)$ le schéma
    des $\omega$ avec $\omega^{\otimes 2} = (-1)^n\delta(q)^{-1}$, un $\mu_2$-torseur. Il
    existe un unique morphisme $\underline{\mathrm{Arf}} \to \underline{\mathrm{Arf}}'$ dont
    le composé avec $\underline{\Delta}_{\{n\}} \to \underline{\mathrm{Arf}}$ soit $F \mapsto
    \omega_F$ ; compatible aux actions de $\mathbf{Z}/2$ et de $\mu_2$, il définit un
    isomorphisme
    \[
      (7) \qquad \underline{\mathrm{Arf}}' \simeq \underline{\mathrm{Arf}}
      \wedge^{(\mathbf{Z}/2)_S} \mu_{2,S},
    \]
    et un isomorphisme $\underline{\mathrm{Arf}} \simeq \underline{\mathrm{Arf}}'$ si $2$
    est inversible.
\end{itemize}
En caractéristique $2$, $\underline{\mathrm{Arf}}$ est le revêtement d'Artin--Schreier qui
porte l'invariant d'Arf, tandis que $\underline{\mathrm{Arf}}'$ est radiciel ; c'est le
$\mathbf{Z}/2$-torseur, non le $\mu_2$-torseur des racines carrées du discriminant, qui
a un sens en toute caractéristique\footnote{Lecture de l'édition. Le nom
« Arf » est sur la page (« vecteurs d'Arf », page 139).}.

\textbf{Application aux groupes orthogonaux} (page 144). Si $N = 2n + 1$, $O(q) \simeq
\mu_2 \times \mathrm{SO}(q)$, $\mathrm{SO}(q)$ lisse à fibres connexes ; $\mathrm{SO}(q)
\simeq O(q)/\mu_2$ opère sur chaque $\underline{\Delta}_H$, qui est un espace homogène.
\textbf{Corollaire} : la projection $O(q) \to \mu_2$ est induite par le déterminant. La
page 145, barrée, commence à prouver que $u \mapsto u(D)$, $\mathrm{SO}(q) \to
\underline{\Delta}_H$, est lisse à fibres connexes, en se ramenant au drapeau complet et
à une base standard $e_0, e_1, \ldots, e_n, e_1^{*}, \ldots, e_n^{*}$.

\pagerange{146}{152}
\subsection*{Scindages isotropes et stabilisateurs (pages 146 à 152)}

\textbf{Proposition} (pages 146 à 148). Soit $(E, q)$ régulière et $F \subset E$ un
sous-Module totalement isotrope facteur direct.
\begin{itemize}
  \item[a)] $E \to E^{\vee} \to F^{\vee}$ est un épimorphisme ; son noyau $F^{\perp}$
    est facteur direct, de formation compatible aux changements de base.
  \item[b)] Pour $G = E/F^{\perp}$, $\varphi$ induit une dualité $F \times G \to
    \mathcal{O}_S$.
  \item[c)] $Q = F^{\perp}/F$ porte une forme régulière $q_Q$, indépendante des choix. Si
    $\theta : G \to E$ scinde $(*)$ $0 \to F^{\perp} \to E \to G \to 0$, la forme $q_\theta(x,
    y) = q(x + \theta y) = q(\theta y) + \langle y, x\rangle$ sur $F \oplus G$ est
    régulière, et $E \simeq (F \oplus \theta G) \perp Q_\theta$, $Q_\theta \simeq Q$.
  \item[d)] Si $S$ est affine, il existe $u : G \to F$ tel que $\theta_u = \theta + u$
    soit un scindage \emph{isotrope} ($q \circ \theta_u = 0$), c'est-à-dire tel que
    $\langle y, u(y)\rangle = -q(\theta y)$ ; les $u$ convenables forment un torseur sous
    $\mathrm{Alt}(G)$, les formes bilinéaires alternées de $G$ ($F \simeq G^{\vee}$). Le
    quotient des scindages isotropes par $\mathrm{Alt}(G)$ s'identifie à celui de tous
    les scindages par $\mathrm{Bil}(G) = \mathrm{Hom}(G, F)$, qui est un torseur sous
    $\mathrm{Hom}(G, Q)$.
  \item[e)] Un scindage isotrope met $q$ sous la forme $\langle x, x'\rangle$ sur $F
    \oplus F^{\vee}$, somme orthogonale avec $Q_\theta$.
\end{itemize}
\textbf{Corollaire.} Un drapeau $0 \subsetneq E_1 \subsetneq \cdots \subsetneq E_\nu =
F$ se complète en $F \subset F^{\perp} \subsetneq E_{\nu - 1}^{\perp} \subsetneq \cdots
\subsetneq E_1^{\perp} \subsetneq E$, et un scindage isotrope identifie $E$ à $F \oplus
F^{\vee} \oplus Q$\footnote{La page termine le drapeau complété par « $F$ » ; c'est
$E$.}.

\textbf{Corollaires 1 à 3.} Si $(E, q)$ et $(E', q')$ sont de même type, $D \in
\Delta_H(E, q)$, $D' \in \Delta_H(E', q')$, $Q \simeq Q'$ isométriquement et $R_i =
E_i/E_{i-1} \simeq R'_i$, et si $S$ est affine, il existe une isométrie $u$ avec $u(D) =
D'$. Localement, il en existe donc toujours, et $\underline{\Delta}_H(E, q)$ est un
espace homogène sous $O(q)$ : $\underline{\Delta}_H \simeq
O(q)/\underline{\mathrm{Stab}}_D(q)$\footnote{La page dit « localement (ét) ». C'est vrai
si $2$ est inversible ou si $Q$ est de rang pair ; si $Q$ est de rang impair en
caractéristique $2$, deux formes régulières de même rang ne sont isomorphes que
localement pour la topologie fidèlement plate (page 127, e)), et c'est elle qu'il faut
prendre.}.

\textbf{Corollaires 4 à 6.} Des isomorphismes $u_i : R_i \simeq R'_i$ et une isométrie
$u_Q : Q \simeq Q'$ étant donnés, $S$ affine, il existe $u$ comme ci-dessus les
induisant, et l'ensemble de ces $u$ est un torseur sous le sous-groupe $O_D(q)^{+}$ des
automorphismes de $(E, q, D)$ qui induisent l'identité sur les $R_i$ et sur $Q$. Soit
$\bar{E} = \bigoplus_i (R_i \oplus R_i^{\vee}) \oplus Q$, avec la forme évidente et le
drapeau $\bar{D}$ ; le produit $\bar{\Sigma} = \prod_i \Sigma_i \times \Sigma_0$ des
scindages des $0 \to E_{i-1} \to E_i \to R_i \to 0$ et des scindages isotropes de $0 \to
E_\nu \to E_\nu^{\perp} \to Q \to 0$ s'identifie au schéma des isomorphismes
$(\bar{E}, \bar{q}, \bar{D}) \to (E, q, D)$ compatibles aux identifications des
gradués, qui est un torseur sous $O_D(q)^{+}$. Comme $\bar{\Sigma}$ est lisse à fibres
connexes : \textbf{Corollaire 6}, $O_D(q)^{+}$ est lisse à fibres connexes, et, si $S$
est affine, isomorphe comme schéma (non comme groupe) à un fibré vectoriel de
dimension relative
\[
  \dim O_D(q)^{+} = \sum_{j < i} d_id_j + \frac{d(d-1)}{2} + d\rho = d^2 + d\rho -
  \sum_i \frac{d_i(d_i+1)}{2},
\]
avec $d_i = \mathrm{rg}\,R_i$, $d = \sum d_i$, $\rho = \mathrm{rg}\,Q$, $N = 2d +
\rho$\footnote{La page décrit ce fibré comme $\check{\mathbb{V}}(\sum_i
\underline{\mathrm{Hom}}(R_i, E_{i-1}) \oplus \underline{\mathrm{Hom}}(\bar{G},
\bar{F}^{\perp})^{!})$, le second terme étant « le noyau de l'épi » $u \mapsto (\bar{y}
\mapsto \bar{q}(\bar{y} + u\bar{y}))$. Cette application n'est pas linéaire ($q(u\bar{y})
= q_Q(\cdot)$ n'est pas nul en général) ; son lieu des zéros est néanmoins, sur une base
affine, isomorphe comme schéma au fibré vectoriel $\mathrm{Hom}(G, Q) \oplus
\mathrm{Alt}(G)$, par d), et la dimension $d\rho + d(d-1)/2$ est celle de la page.}.
D'où, avec $\dim(\prod \mathrm{GL}(R_i) \times O(q_Q)) = \sum d_i^2 + \rho(\rho -
1)/2$ et $\dim O(q) = N(N-1)/2$ :
\[
  \dim O_D(q) = d^2 + d\rho + \frac{\rho(\rho - 1)}{2} + \sum_i \frac{d_i(d_i - 1)}{2},
  \qquad
  \dim \underline{\Delta}_H(q) = d^2 + d\rho - d - \sum_i \frac{d_i(d_i - 1)}{2}.
\]
Par exemple $\dim \underline{\Delta}_{\{d\}} = \frac{d(d+1)}{2} + d(\rho - 1)$, et, pour $N
= 2n$, $d = n$, $\rho = 0$ : $\frac{n(n-1)}{2}$, la dimension de la grassmannienne
orthogonale\footnote{La page 153 a « $+\,d(\rho - 1)$ » sur un signe surchargé ; la
formule générale le confirme.}.

\pagerange{153}{157}
\subsection*{$\mathrm{SO}(q)$ lisse à fibres connexes ; l'invariant de Dickson (pages 153 à 157)}

\textbf{Rang impair.} Pour $D \in \Delta_H$, le stabilisateur de $D$ dans
$\mathrm{SO}(q)$ est $\mathrm{SO}_D(q) = O_D(q) \cap \mathrm{SL}(E)$. Le déterminant sur
$O_D(q)$ se factorise par $\alpha : O_D(q) \to \prod_i \mathrm{GL}(R_i) \times O(q_Q)$ et
$\det_Q$ (la partie $F \oplus F^{\vee}$ contribue $\det g \cdot \det g^{-1} = 1$), d'où
$\mathrm{SO}_D(q) = \alpha^{-1}(\prod \mathrm{GL}(R_i) \times \mathrm{SO}(q_Q))$ et
(Proposition de la page 154) un diagramme d'extensions
\[
  1 \to O_D(q)^{+} \to \mathrm{SO}_D(q) \to \prod_i \mathrm{GL}(R_i) \times
  \mathrm{SO}(q_Q) \to 1 ,
\]
où $O_D(q)^{+} = \mathrm{SO}_D(q)^{+}$ : $\mathrm{SO}_D(q)$ est lisse à fibres connexes.

\textbf{Rang pair, première tentative} (pages 154 et 155). La page énonce que
$\mathrm{SO}(q) := O(q) \cap \mathrm{SL}(E)$ est lisse à fibres connexes : pour $F$
totalement isotrope de rang $n$, $Q = 0$\footnote{La page écrit « $F^{\perp}/F = Q$ de rg
$1$ », ce qui est le cas $N = 2n + 1$ ; pour $N = 2n$, $Q = 0$.}, $\mathrm{SO}_D(q)$ est
extension de $\mathrm{GL}(F)$ par un groupe lisse à fibres connexes, et
$\mathrm{SO}(q)/\mathrm{SO}_D(q) \simeq \underline{\Delta}_{\{n\}}$ ; mais en
caractéristique $2$, $\det$ est trivial sur $O(q)$ et $O(q) \cap \mathrm{SL}(E) = O(q)$
n'est pas connexe, et $\underline{\Delta}_{\{n\}}$ non plus. \textbf{Corollaire} : en
rang impair, les $\mathrm{SO}_D(q) = B_D(q)$ sont lisses à fibres connexes et $O_D(q)
\simeq B_D(q) \times \mu_2$.

\textbf{Rang pair, la bonne définition} (pages 155 à 157). « On ne va pas définir
$\mathrm{SO}(q)$ comme $O(q) \cap \mathrm{SL}(E)$, à cause de car.\ $2$. » On fait opérer
$O(q)$ sur $\underline{\mathrm{Arf}}(q)$, d'où un homomorphisme
\[
  \varepsilon_q : O(q) \to (\mathbf{Z}/2)_S, \qquad \mathrm{SO}(q) :=
  \mathrm{Ker}\,\varepsilon_q, \qquad 1 \to \mathrm{SO}(q) \to O(q) \to (\mathbf{Z}/2)_S \to
  1 ;
\]
si $\eta : \mathbf{Z}/2 \to \mu_2$ est le morphisme canonique, $\eta \circ \varepsilon_q =
\det$, de sorte que $\mathrm{SO}(q) = O(q) \cap \mathrm{SL}(E)$ si $2$ est inversible
(marge). C'est l'invariant de Dickson\footnote{Le nom n'est pas sur la page.}.
\textbf{Proposition} : $\varepsilon_q|O_D(q)$ est le composé $O_D(q) \to O(q_Q)
\xrightarrow{\varepsilon_{q_Q}} \mathbf{Z}/2$, donc $\mathrm{SO}_D(q)$ est l'image
réciproque de $\mathrm{SO}(q_Q)$, extension de $\prod \mathrm{GL}(R_i) \times
\mathrm{SO}(q_Q)$ par $O_D(q)^{+}$, lisse à fibres connexes.

\textbf{Théorème.} Si $(E, q)$ est régulière de rang $2n$, $\mathrm{SO}(q)$ est lisse à
fibres connexes. On peut supposer qu'il existe $F \in \Delta_{\{n\}}(q)$, d'où une section
$\omega$ de $\underline{\mathrm{Arf}}(q)$ ; pour $D = \{F\}$, $Q = 0$, $\mathrm{SO}_D(q)$
est lisse à fibres connexes, et $\mathrm{SO}(q)/\mathrm{SO}_D(q) \simeq
\underline{\Delta}^{\omega}_{\{n\}}$ aussi ; donc $\mathrm{SO}(q)$ l'est, cqfd.
\textbf{Corollaire} : les $\mathrm{SO}_D(q) = B_D(q)$ — les sous-groupes paraboliques,
de Borel pour un drapeau complet — sont lisses à fibres connexes.

En marge de la page 155, une \textbf{Proposition} sans démonstration :
$\mathrm{SO}(q)$ opère transitivement sur $\underline{\Delta}_H(q)$ si $N$ est impair
ou $d = \max H < n$ ; si $N = 2n$ et $d = n$, et si $\omega$ est une section de
$\underline{\mathrm{Arf}}(q)$, il opère transitivement sur la partie
$\underline{\Delta}^{\omega}_H(q)$ au-dessus de $\omega$\footnote{La marge porte aussi
« et aussi si $S \neq \varnothing$ », qu'on ne sait pas rattacher.}. Le dossier
s'arrête sur le corollaire de la page 157 ; le reste du feuillet est blanc.

\end{document}
